%<<>>==<<>><<>>==<<>><<>>==<<>><<>>==<<>><<>>==<<>><<>>==<<>>
\documentclass[12pt]{article}
%<<>>==<<>><<>>==<<>><<>>==<<>><<>>==<<>><<>>==<<>><<>>==<<>>
\usepackage[utf8]{inputenc}
\usepackage[T1]{fontenc}
\usepackage[spanish,mexico]{babel}
\usepackage{amsmath,amssymb,amsthm,mathtools}
\usepackage{geometry}
\usepackage{enumitem}
\usepackage{booktabs}
\usepackage{array}
\usepackage{xcolor}
\usepackage{hyperref}
\usepackage{fancyhdr}
\usepackage{microtype}
\usepackage{longtable}
\usepackage{multicol}
%<<>>==<<>><<>>==<<>><<>>==<<>><<>>==<<>><<>>==<<>><<>>==<<>>
\geometry{letterpaper,margin=2.35cm}
\setlength{\parindent}{0pt}
\setlength{\parskip}{0.55em}
\setlength{\headheight}{15pt}
\setlist[itemize]{leftmargin=1.6em}
\setlist[enumerate]{leftmargin=1.9em}
%<<>>==<<>><<>>==<<>><<>>==<<>><<>>==<<>><<>>==<<>><<>>==<<>>
\definecolor{azul}{RGB}{22,72,120}
\hypersetup{colorlinks=true,linkcolor=azul,urlcolor=azul}
%<<>>==<<>><<>>==<<>><<>>==<<>><<>>==<<>><<>>==<<>><<>>==<<>>
\pagestyle{fancy}
\fancyhf{}
\lhead{Tecnológico Nacional de México}
\rhead{Álgebra Lineal}
\cfoot{\thepage}
%<<>>==<<>><<>>==<<>><<>>==<<>><<>>==<<>><<>>==<<>><<>>==<<>>
\newtheorem{definicion}{Definición}[section]
\newtheorem{teorema}[definicion]{Teorema}
\newtheorem{proposicion}[definicion]{Proposición}
\newtheorem{corolario}[definicion]{Corolario}
\newtheorem{Nota}[definicion]{Nota}
\theoremstyle{remark}
\newtheorem{observacion}[definicion]{Observación}
\newtheorem{ejemplo}[definicion]{Ejemplo}
\newtheorem{Ejem}[definicion]{Ejemplo}
\newtheorem{Ejer}[definicion]{Ejercicio}
%<<>>==<<>><<>>==<<>><<>>==<<>><<>>==<<>><<>>==<<>><<>>==<<>>
% Niveles de dificultad para el banco de ejercicios
\newcommand{\basico}{\textbf{[Básico]}}
\newcommand{\intermedio}{\textbf{[Intermedio]}}
\newcommand{\avanzado}{\textbf{[Avanzado]}}
\newcommand{\reto}{\textbf{[Reto]}}
\newcommand{\R}{\mathbb{R}}
\newcommand{\rank}{\operatorname{rang}}
\newcommand{\RREF}{\operatorname{RREF}}
\newcommand{\sistema}[1]{%
\[
\left\{
\begin{aligned}
#1
\end{aligned}
\right.
\]
}
%<<>>==<<>><<>>==<<>><<>>==<<>><<>>==<<>><<>>==<<>><<>>==<<>>
\title{\textbf{Notas para el curso de Álgebra Lineal}\\[0.35cm]
\Large Unidad 2: Sistemas de ecuaciones lineales}
\author{Carlos Ernesto Martínez\\Tecnológico Nacional de México}
\date{Agosto--Diciembre 2026}
%<<>>==<<>><<>>==<<>><<>>==<<>><<>>==<<>><<>>==<<>><<>>==<<>>
\begin{document}
\maketitle
\tableofcontents
\newpage
%<<>>==<<>><<>>==<<>><<>>==<<>><<>>==<<>><<>>==<<>><<>>==<<>>
\section*{Propósito de la unidad}
%<<>>==<<>><<>>==<<>><<>>==<<>><<>>==<<>><<>>==<<>><<>>==<<>>
En esta unidad estudiaremos los \emph{sistemas de ecuaciones lineales}. Nos interesa comprender qué significa resolver simultáneamente varias ecuaciones, cuándo un sistema tiene solución, cuántas soluciones puede tener, y cómo interpretar geométricamente esas posibilidades.

La idea central será pasar progresivamente de la escritura usual de un sistema,
\[
\begin{cases}
a_{11}x_1+\cdots+a_{1n}x_n=b_1,\\
\vdots\\
a_{m1}x_1+\cdots+a_{mn}x_n=b_m,
\end{cases}
\]
a una representación matricial compacta
\[
A\mathbf{x}=\mathbf{b}.
\]
La matriz no sustituye al sistema: conserva de manera organizada la información de sus coeficientes y permite realizar de forma eficiente las mismas operaciones algebraicas que ya conocemos.

%<<>>==<<>><<>>==<<>><<>>==<<>><<>>==<<>><<>>==<<>><<>>==<<>>
\section{Sistemas de Ecuaciones lineales}
%<<>>==<<>><<>>==<<>><<>>==<<>><<>>==<<>><<>>==<<>><<>>==<<>>
\begin{definicion}
Una \textbf{ecuación lineal} en las incógnitas $x_1,\ldots,x_n$ es una ecuación de la forma
\[
a_1x_1+a_2x_2+\cdots+a_nx_n=b,
\]
donde $a_1,\ldots,a_n,b\in\R$ son constantes conocidas.
\end{definicion}


\begin{ejemplo}
Las expresiones
\[
2x-3y=7,\qquad x_1+4x_2-2x_3=5,\qquad 3x+2y-z=0
\]
son ecuaciones lineales.
\end{ejemplo}

\begin{ejemplo}
Las ecuaciones
\[
x^2+y=3,\qquad xy+z=1,\qquad \sin x+y=0
\]
no son lineales.
\end{ejemplo}

\begin{definicion}
Una \textbf{solución} de una ecuación lineal es un conjunto de valores de las incógnitas que hace verdadera
la igualdad.
\end{definicion}

Por ejemplo, $(x,y)=(2,1)$ es solución de $2x-y=3$, pues
\[
2(2)-1=3.
\]
Sin embargo, una sola ecuación con varias incógnitas suele tener muchas soluciones. El problema se vuelve
más interesante cuando exigimos que los mismos valores satisfagan simultáneamente varias ecuaciones.
%<<>>==<<>><<>>==<<>><<>>==<<>><<>>==<<>><<>>==<<>><<>>==<<>>
%\section{Sistemas de ecuaciones lineales}
%<<>>==<<>><<>>==<<>><<>>==<<>><<>>==<<>><<>>==<<>><<>>==<<>>
\begin{definicion}
Un \textbf{sistema de $m$ ecuaciones lineales con $n$ incógnitas} es un conjunto de ecuaciones
\[
\begin{cases}
a_{11}x_1+a_{12}x_2+\cdots+a_{1n}x_n=b_1,\\
a_{21}x_1+a_{22}x_2+\cdots+a_{2n}x_n=b_2,\\
\hspace{2.5cm}\vdots\\
a_{m1}x_1+a_{m2}x_2+\cdots+a_{mn}x_n=b_m.
\end{cases}
\]
\end{definicion}

Una solución del sistema debe satisfacer \emph{todas} las ecuaciones al mismo tiempo.

\begin{ejemplo}
Consideremos
\[
\begin{cases}
x+y=5,\\
2x-y=1.
\end{cases}
\]
Sumando ambas ecuaciones,
\[
3x=6,
\]
de donde $x=2$. Sustituyendo en la primera ecuación,
\[
2+y=5,
\]
y por tanto $y=3$. La solución es
\[
(x,y)=(2,3).
\]
La verificación es inmediata:
\[
2+3=5,\qquad 2(2)-3=1.
\]
\end{ejemplo}
%<<>>==<<>><<>>==<<>><<>>==<<>><<>>==<<>><<>>==<<>><<>>==<<>>
%\subsection{¿Qué significa resolver un sistema?}
%<<>>==<<>><<>>==<<>><<>>==<<>><<>>==<<>><<>>==<<>><<>>==<<>>
Resolver un sistema significa describir \emph{todo su conjunto solución}. No basta encontrar una pareja que funcione si existen otras. Tampoco debemos concluir que ``no se puede resolver'' sólo porque un procedimiento produce una variable libre: precisamente esa variable libre puede indicar la existencia de infinitas soluciones.
%<<>>==<<>><<>>==<<>><<>>==<<>><<>>==<<>><<>>==<<>><<>>==<<>>
\section*{Interpretación geométrica}
%<<>>==<<>><<>>==<<>><<>>==<<>><<>>==<<>><<>>==<<>><<>>==<<>>
La geometría permite comprender, antes de desarrollar los algoritmos, por qué aparecen distintos tipos de solución.
%<<>>==<<>><<>>==<<>><<>>==<<>><<>>==<<>><<>>==<<>><<>>==<<>>
%\subsection{Dos ecuaciones con dos incógnitas}
%<<>>==<<>><<>>==<<>><<>>==<<>><<>>==<<>><<>>==<<>><<>>==<<>>
Una ecuación
\[
ax+by=c
\]
representa, salvo casos degenerados, una recta en $\R^2$. Resolver dos ecuaciones simultáneamente equivale a buscar los puntos comunes de dos rectas.

Existen tres posibilidades fundamentales.
%=============================================
\subsubsection*{Rectas que se cortan}
%=============================================
\[
\begin{cases}
x+y=3,\\
x-y=1.
\end{cases}
\]
Las rectas se intersectan en un único punto. El sistema tiene una solución.
%=============================================
\subsubsection*{Rectas paralelas distintas}
%=============================================
\[
\begin{cases}
x+y=2,\\
2x+2y=5.
\end{cases}
\]
La segunda ecuación tendría el mismo lado izquierdo que el doble de la primera, pero un término independiente diferente. Las rectas son paralelas y el sistema no tiene solución.
%=============================================
\subsubsection*{La misma recta}
%=============================================
\[
\begin{cases}
x+y=2,\\
2x+2y=4.
\end{cases}
\]
La segunda ecuación es exactamente el doble de la primera. Ambas describen la misma recta, por lo que cada punto de ella es solución: existen infinitas soluciones.
%=============================================
\subsection*{Tres incógnitas}
%=============================================
En $\R^3$, una ecuación lineal
\[
ax+by+cz=d
\]
representa normalmente un plano. Resolver un sistema significa buscar la intersección común de los planos. Pueden cortarse en un punto, en una recta, coincidir parcialmente o no tener ningún punto común. Esta interpretación anticipa un hecho fundamental.

\begin{proposicion}
Un sistema lineal puede tener:
\begin{itemize}
\begin{multicols}{3}
\item una única solución;
\item ninguna solución;
\item infinitas soluciones.
\end{multicols}
\end{itemize}
No puede tener exactamente dos, tres o cualquier otro número finito mayor que uno de soluciones.
\end{proposicion}

%<<>>==<<>><<>>==<<>><<>>==<<>><<>>==<<>><<>>==<<>><<>>==<<>>
\section{Representación matricial}
%<<>>==<<>><<>>==<<>><<>>==<<>><<>>==<<>><<>>==<<>><<>>==<<>>
Consideremos
\[
\begin{cases}
2x-y+3z=7,\\
x+4y-z=2,\\
3x+2y+2z=5.
\end{cases}
\]
Separaremos tres objetos:
\[
A=
\begin{pmatrix}
2&-1&3\\
1&4&-1\\
3&2&2
\end{pmatrix},
\qquad
\mathbf{x}=
\begin{pmatrix}x\\y\\z\end{pmatrix},
\qquad
\mathbf{b}=
\begin{pmatrix}7\\2\\5\end{pmatrix}.
\]
Entonces el sistema se escribe
\[
A\mathbf{x}=\mathbf{b}.
\]

\begin{definicion}
$A$ se denomina \textbf{matriz de coeficientes}; $\mathbf{x}$ es el vector de incógnitas y
$\mathbf{b}$ el vector de términos independientes.
\end{definicion}

\begin{definicion}
La \textbf{matriz aumentada} del sistema es
\[
[A\mid\mathbf b]=
\left[
\begin{array}{ccc|c}
2&-1&3&7\\
1&4&-1&2\\
3&2&2&5
\end{array}
\right].
\]
\end{definicion}

La barra vertical es sólo una ayuda visual: separa los coeficientes de las incógnitas del término
independiente.
%<<>>==<<>><<>>==<<>><<>>==<<>><<>>==<<>><<>>==<<>><<>>==<<>>
\subsection*{Sistemas homogéneos y no homogéneos}
%<<>>==<<>><<>>==<<>><<>>==<<>><<>>==<<>><<>>==<<>><<>>==<<>>
\begin{definicion}
Un sistema es \textbf{homogéneo} si todos sus términos independientes son cero:
\[
A\mathbf{x}=\mathbf{0}.
\]
Es \textbf{no homogéneo} si
\[
A\mathbf{x}=\mathbf{b},\qquad \mathbf b\neq\mathbf0.
\]
\end{definicion}

Todo sistema homogéneo posee al menos una solución:
\[
\mathbf{x}=\mathbf0.
\]
Ésta se denomina \textbf{solución trivial}.

\begin{ejemplo}
\[
\begin{cases}
x+2y-z=0,\\
2x+4y-2z=0.
\end{cases}
\]
La segunda ecuación es el doble de la primera. Tomemos
\[
y=s,\qquad z=t.
\]
Entonces
\[
x=-2s+t,
\]
y todas las soluciones son
\[
\begin{pmatrix}x\\y\\z\end{pmatrix}
=
\begin{pmatrix}-2s+t\\s\\t\end{pmatrix}
=
s\begin{pmatrix}-2\\1\\0\end{pmatrix}
+t\begin{pmatrix}1\\0\\1\end{pmatrix}.
\]
Además de la solución trivial $s=t=0$, existen infinitas soluciones no triviales.
\end{ejemplo}

%<<>>==<<>><<>>==<<>><<>>==<<>><<>>==<<>><<>>==<<>><<>>==<<>>
\subsection*{Sistemas consistenes e inconsistentes}
%<<>>==<<>><<>>==<<>><<>>==<<>><<>>==<<>><<>>==<<>><<>>==<<>>

\begin{definicion}
Un sistema es \textbf{compatible} si tiene al menos una solución. Es \textbf{incompatible} si no tiene solución.
\end{definicion}

Un sistema compatible puede ser:
\begin{itemize}
\item \textbf{compatible determinado}: tiene una única solución;
\item \textbf{compatible indeterminado}: tiene infinitas soluciones.
\end{itemize}

%<<>>==<<>><<>>==<<>><<>>==<<>><<>>==<<>><<>>==<<>><<>>==<<>>
\subsection*{De la eliminación algebraica a las operaciones por renglones}
%<<>>==<<>><<>>==<<>><<>>==<<>><<>>==<<>><<>>==<<>><<>>==<<>>
Antes de introducir un algoritmo matricial, recordemos la eliminación usual:
\[
\begin{cases}
x+2y=5,\\
3x-y=4.
\end{cases}
\]
Para eliminar $x$, reemplazamos la segunda ecuación por
\[
E_2-3E_1.
\]
Obtenemos
\[
-7y=-11.
\]
Lo importante es observar que no hemos cambiado el conjunto solución: sólo reemplazamos una ecuación por otra consecuencia algebraicamente equivalente.

La matriz aumentada es
\[
\left[
\begin{array}{cc|c}
1&2&5\\
3&-1&4
\end{array}
\right].
\]
La misma operación se escribe
\[
R_2\leftarrow R_2-3R_1,
\]
y produce
\[
\left[
\begin{array}{cc|c}
1&2&5\\
0&-7&-11
\end{array}
\right].
\]
%<<>>==<<>><<>>==<<>><<>>==<<>><<>>==<<>><<>>==<<>><<>>==<<>>
\subsection*{Operaciones elementales}
%<<>>==<<>><<>>==<<>><<>>==<<>><<>>==<<>><<>>==<<>><<>>==<<>>

\begin{definicion}
Las \textbf{operaciones elementales por renglones} son:
\begin{enumerate}
\item intercambiar dos renglones:
\[
R_i\leftrightarrow R_j;
\]
\item multiplicar un renglón por un escalar no nulo:
\[
R_i\leftarrow cR_i,\qquad c\neq0;
\]
\item sumar a un renglón un múltiplo de otro:
\[
R_i\leftarrow R_i+cR_j.
\]
\end{enumerate}
\end{definicion}

\begin{teorema}
Las operaciones elementales por renglones preservan el conjunto solución del sistema.
\end{teorema}

\begin{Nota}
Intercambiar ecuaciones no modifica la exigencia de que todas se satisfagan. Multiplicar una ecuación
por un número no nulo produce una ecuación equivalente. Finalmente, sustituir una ecuación por ella
misma más un múltiplo de otra es reversible:
\[
R_i\leftarrow R_i+cR_j
\quad\Longleftrightarrow\quad
R_i\leftarrow R_i-cR_j.
\]
Por tanto, cada operación transforma el sistema en otro equivalente.
\end{Nota}
%<<>>==<<>><<>>==<<>><<>>==<<>><<>>==<<>><<>>==<<>><<>>==<<>>
\subsection*{Forma escalonada}
%<<>>==<<>><<>>==<<>><<>>==<<>><<>>==<<>><<>>==<<>><<>>==<<>>
\begin{definicion}
Una matriz está en \textbf{forma escalonada por renglones} si:
\begin{itemize}
\item todos los renglones nulos están debajo de los no nulos;
\item el primer elemento no nulo de cada renglón aparece a la derecha del primer elemento no nulo
del renglón anterior;
\item debajo de cada pivote hay ceros.
\end{itemize}
\end{definicion}

Por ejemplo,
\[
\begin{pmatrix}
1&2&-1&3\\
0&3&4&2\\
0&0&5&-1\\
0&0&0&0
\end{pmatrix}
\]
está en forma escalonada.

\begin{definicion}
La primera entrada no nula de un renglón no nulo se denomina \textbf{pivote}. La columna que contiene
un pivote se denomina \textbf{columna pivote}.
\end{definicion}
%<<>>==<<>><<>>==<<>><<>>==<<>><<>>==<<>><<>>==<<>><<>>==<<>>
\section{Eliminación de Gauss}
%<<>>==<<>><<>>==<<>><<>>==<<>><<>>==<<>><<>>==<<>><<>>==<<>>
El método de Gauss consiste en aplicar operaciones elementales hasta llevar la matriz aumentada a forma escalonada. Después se utiliza sustitución hacia atrás.

\subsection*{Ejemplo con solución única}

Resolvamos
\[
\begin{cases}
x+y+z=6,\\
2x-y+z=3,\\
x+2y-z=2.
\end{cases}
\]
La matriz aumentada es
\[
\left[
\begin{array}{ccc|c}
1&1&1&6\\
2&-1&1&3\\
1&2&-1&2
\end{array}
\right].
\]
Aplicamos
\[
R_2\leftarrow R_2-2R_1,\qquad
R_3\leftarrow R_3-R_1,
\]
para obtener
\[
\left[
\begin{array}{ccc|c}
1&1&1&6\\
0&-3&-1&-9\\
0&1&-2&-4
\end{array}
\right].
\]
Conviene intercambiar los últimos dos renglones:
\[
R_2\leftrightarrow R_3,
\]
de modo que
\[
\left[
\begin{array}{ccc|c}
1&1&1&6\\
0&1&-2&-4\\
0&-3&-1&-9
\end{array}
\right].
\]
Ahora
\[
R_3\leftarrow R_3+3R_2,
\]
y obtenemos
\[
\left[
\begin{array}{ccc|c}
1&1&1&6\\
0&1&-2&-4\\
0&0&-7&-21
\end{array}
\right].
\]
La última ecuación da $z=3$. Luego
\[
y-2(3)=-4\quad\Longrightarrow\quad y=2,
\]
y finalmente
\[
x+2+3=6\quad\Longrightarrow\quad x=1.
\]
Así,
\[
\boxed{(x,y,z)=(1,2,3)}.
\]
%<<>>==<<>><<>>==<<>><<>>==<<>><<>>==<<>><<>>==<<>><<>>==<<>>
\subsection*{Ejemplo inconsistente}
%<<>>==<<>><<>>==<<>><<>>==<<>><<>>==<<>><<>>==<<>><<>>==<<>>
\[
\begin{cases}
x+y=2,\\
2x+2y=5.
\end{cases}
\]
Tenemos
\[
\left[
\begin{array}{cc|c}
1&1&2\\
2&2&5
\end{array}
\right]
\xrightarrow{R_2-2R_1}
\left[
\begin{array}{cc|c}
1&1&2\\
0&0&1
\end{array}
\right].
\]
El último renglón representa
\[
0x+0y=1,
\]
es decir,
\[
0=1,
\]
una contradicción. Por tanto, el sistema es inconsistente.
%<<>>==<<>><<>>==<<>><<>>==<<>><<>>==<<>><<>>==<<>><<>>==<<>>
\subsection*{Ejemplo con infinitas soluciones}
%<<>>==<<>><<>>==<<>><<>>==<<>><<>>==<<>><<>>==<<>><<>>==<<>>
\[
\begin{cases}
x+y+z=2,\\
2x+2y+2z=4.
\end{cases}
\]
La reducción produce
\[
\left[
\begin{array}{ccc|c}
1&1&1&2\\
0&0&0&0
\end{array}
\right].
\]
Sólo existe una ecuación independiente para tres incógnitas. Tomemos
\[
y=s,\qquad z=t.
\]
Entonces
\[
x=2-s-t.
\]
Por tanto,
\[
\boxed{(x,y,z)=(2-s-t,s,t),\qquad s,t\in\R.}
\]
%<<>>==<<>><<>>==<<>><<>>==<<>><<>>==<<>><<>>==<<>><<>>==<<>>
\subsection*{Variables básicas y variables libres}
%<<>>==<<>><<>>==<<>><<>>==<<>><<>>==<<>><<>>==<<>><<>>==<<>>
En una matriz escalonada, las variables correspondientes a columnas pivote se denominan \textbf{variables básicas}. Las restantes son \textbf{variables libres}.

En
\[
\left[
\begin{array}{cccc|c}
1&2&0&-1&3\\
0&0&1&4&2
\end{array}
\right],
\]
las columnas pivote son la primera y la tercera. Por tanto,
\[
x_1,\ x_3
\]
son básicas, mientras que
\[
x_2,\ x_4
\]
son libres.

Tomando
\[
x_2=s,\qquad x_4=t,
\]
obtenemos
\[
x_1=3-2s+t,\qquad x_3=2-4t.
\]
La solución completa es
\[
\mathbf{x}=
\begin{pmatrix}3\\0\\2\\0\end{pmatrix}
+s\begin{pmatrix}-2\\1\\0\\0\end{pmatrix}
+t\begin{pmatrix}1\\0\\-4\\1\end{pmatrix}.
\]
%<<>>==<<>><<>>==<<>><<>>==<<>><<>>==<<>><<>>==<<>><<>>==<<>>
\section*{Forma escalonada reducida}
%<<>>==<<>><<>>==<<>><<>>==<<>><<>>==<<>><<>>==<<>><<>>==<<>>

\begin{definicion}
Una matriz está en \textbf{forma escalonada reducida por renglones} si está en forma escalonada y, además:
\begin{itemize}
\item cada pivote es igual a $1$;
\item cada pivote es el único elemento no nulo de su columna.
\end{itemize}
\end{definicion}

Por ejemplo,
\[
\begin{pmatrix}
1&0&2&0&3\\
0&1&-1&0&4\\
0&0&0&1&-2
\end{pmatrix}
\]
está en forma escalonada reducida.

\begin{teorema}
Toda matriz es equivalente por renglones a una única matriz escalonada reducida.
\end{teorema}

Esta unicidad es importante: podemos realizar operaciones distintas durante el proceso, pero si continuamos hasta la forma reducida, el resultado final será el mismo.

%<<>>==<<>><<>>==<<>><<>>==<<>><<>>==<<>><<>>==<<>><<>>==<<>>
\section{Método de Gauss--Jordan}
%<<>>==<<>><<>>==<<>><<>>==<<>><<>>==<<>><<>>==<<>><<>>==<<>>

El método de Gauss--Jordan es una continuación sistemática de la eliminación de Gauss. La meta ya no es detenernos cuando la matriz queda escalonada, sino continuar hasta obtener su \textbf{forma escalonada reducida por renglones} (RREF). En una columna pivote deben ocurrir dos cosas: el pivote debe ser \(1\) y todas las demás entradas de esa columna deben ser \(0\).
%<<>>==<<>><<>>==<<>><<>>==<<>><<>>==<<>><<>>==<<>><<>>==<<>>
\subsection*{Algoritmo de trabajo}
%<<>>==<<>><<>>==<<>><<>>==<<>><<>>==<<>><<>>==<<>><<>>==<<>>
Para resolver un sistema por Gauss--Jordan:

\begin{itemize}
\item Escribimos la matriz aumentada.
\item Elegimos la primera columna disponible que pueda contener un pivote.
\item Si el candidato a pivote es cero, intercambiamos renglones.
\item Normalizamos el renglón pivote para que el pivote sea \(1\).
\item Utilizamos ese renglón para producir ceros \emph{debajo y encima} del pivote.
\item Pasamos a la siguiente columna y repetimos el proceso.
\item Al llegar a RREF, interpretamos la matriz: solución única, variables libres o contradicción.
\item Siempre que sea razonable, comprobamos el resultado en el sistema original.
\end{itemize}

\begin{Nota}
En sistemas grandes es especialmente importante escribir cada operación elemental. Una operación no documentada hace muy difícil localizar un error aritmético varias etapas después.
\end{Nota}

%<<>>==<<>><<>>==<<>><<>>==<<>><<>>==<<>><<>>==<<>><<>>==<<>>
%\subsection*{Ejemplo guiado 1: sistema $2\times2$}
%<<>>==<<>><<>>==<<>><<>>==<<>><<>>==<<>><<>>==<<>><<>>==<<>>

\begin{Ejem}
Consideremos el sistema
\[
\begin{aligned}
x_{1} + 2 x_{2} &= 5\\
2 x_{1} - x_{2} &= 5
\end{aligned}
\]
La matriz aumentada correspondiente es
\[
\left[\begin{array}{cc|c}
1 & 2 & 5\\
2 & -1 & 5
\end{array}\right]
\]
Aplicaremos Gauss--Jordan de manera sistemática. En cada etapa indicaremos qué entrada queremos convertir en pivote y qué entradas deseamos anular.


Usamos el pivote de la columna 1 para anular la entrada del renglón 2.
\[
R_{2}\leftarrow R_{2}-2R_{1}
\]
\[
\left[\begin{array}{cc|c}
1 & 2 & 5\\
0 & -5 & -5
\end{array}\right]
\]
Normalizamos el pivote de la columna 2 para convertirlo en 1.
\[
R_{2}\leftarrow \frac{1}{-5}R_{2}
\]
\[
\left[\begin{array}{cc|c}
1 & 2 & 5\\
0 & 1 & 1
\end{array}\right]
\]
Usamos el pivote de la columna 2 para anular la entrada del renglón 1.
\[
R_{1}\leftarrow R_{1}-2R_{2}
\]
\[
\left[\begin{array}{cc|c}
1 & 0 & 3\\
0 & 1 & 1
\end{array}\right]
\]



Al finalizar obtenemos la forma escalonada reducida. La solución puede leerse directamente de la última columna.
\[
\boxed{x_{1}=3,\quad x_{2}=1}
\]
\end{Ejem}
%==========================================
\begin{Ejem}
Consideremos el sistema

\[
\begin{aligned}
x_{1} + x_{2} + x_{3} &= 6\\
2 x_{1} - x_{2} + x_{3} &= 3\\
x_{1} + 2 x_{2} - x_{3} &= 2
\end{aligned}
\]
La matriz aumentada correspondiente es
\[
\left[\begin{array}{ccc|c}
1 & 1 & 1 & 6\\
2 & -1 & 1 & 3\\
1 & 2 & -1 & 2
\end{array}\right]
\]
Usamos el pivote de la columna 1 para anular la entrada del renglón 2.
\[
R_{2}\leftarrow R_{2}-2R_{1}
\]
\[
\left[\begin{array}{ccc|c}
1 & 1 & 1 & 6\\
0 & -3 & -1 & -9\\
1 & 2 & -1 & 2
\end{array}\right]
\]
 Usamos el pivote de la columna 1 para anular la entrada del renglón 3.
\[
R_{3}\leftarrow R_{3}-R_{1}
\]
\[
\left[\begin{array}{ccc|c}
1 & 1 & 1 & 6\\
0 & -3 & -1 & -9\\
0 & 1 & -2 & -4
\end{array}\right]
\]
 Normalizamos el pivote de la columna 2 para convertirlo en 1.
\[
R_{2}\leftarrow \frac{1}{-3}R_{2}
\]
\[
\left[\begin{array}{ccc|c}
1 & 1 & 1 & 6\\
0 & 1 & \frac{1}{3} & 3\\
0 & 1 & -2 & -4
\end{array}\right]
\]
 Usamos el pivote de la columna 2 para anular la entrada del renglón 1.
\[
R_{1}\leftarrow R_{1}-R_{2}
\]
\[
\left[\begin{array}{ccc|c}
1 & 0 & \frac{2}{3} & 3\\
0 & 1 & \frac{1}{3} & 3\\
0 & 1 & -2 & -4
\end{array}\right]
\]
 Usamos el pivote de la columna 2 para anular la entrada del renglón 3.
\[
R_{3}\leftarrow R_{3}-R_{2}
\]
\[
\left[\begin{array}{ccc|c}
1 & 0 & \frac{2}{3} & 3\\
0 & 1 & \frac{1}{3} & 3\\
0 & 0 & - \frac{7}{3} & -7
\end{array}\right]
\]
Normalizamos el pivote de la columna 3 para convertirlo en 1.
\[
R_{3}\leftarrow \frac{1}{- \frac{7}{3}}R_{3}
\]
\[
\left[\begin{array}{ccc|c}
1 & 0 & \frac{2}{3} & 3\\
0 & 1 & \frac{1}{3} & 3\\
0 & 0 & 1 & 3
\end{array}\right]
\]
Usamos el pivote de la columna 3 para anular la entrada del renglón 1.
\[
R_{1}\leftarrow R_{1}-\frac{2}{3}R_{3}
\]
\[
\left[\begin{array}{ccc|c}
1 & 0 & 0 & 1\\
0 & 1 & \frac{1}{3} & 3\\
0 & 0 & 1 & 3
\end{array}\right]
\]
Usamos el pivote de la columna 3 para anular la entrada del renglón 2.
\[
R_{2}\leftarrow R_{2}-\frac{1}{3}R_{3}
\]
\[
\left[\begin{array}{ccc|c}
1 & 0 & 0 & 1\\
0 & 1 & 0 & 2\\
0 & 0 & 1 & 3
\end{array}\right]
\]
Al finalizar obtenemos la forma escalonada reducida. La solución puede leerse directamente de la última columna.
\[
\boxed{x_{1}=1,\quad x_{2}=2,\quad x_{3}=3}
\]

\end{Ejem}
%==========================================
\begin{Ejem}

Partimos del sistema
\[
\begin{aligned}
x_{2} + x_{3} &= 2\\
x_{1} - x_{2} + 2 x_{3} &= 9\\
2 x_{1} + x_{2} - x_{3} &= 0
\end{aligned}
\]
La matriz aumentada correspondiente es
\[
\left[\begin{array}{ccc|c}
0 & 1 & 1 & 2\\
1 & -1 & 2 & 9\\
2 & 1 & -1 & 0
\end{array}\right]
\]
Aplicaremos Gauss--Jordan de manera sistemática. En cada etapa indicaremos qué entrada queremos convertir en pivote y qué entradas deseamos anular.


Intercambiamos renglones para colocar un elemento no nulo en la posición pivote de la columna 1.
\[
R_{1}\leftrightarrow R_{2}
\]
\[
\left[\begin{array}{ccc|c}
1 & -1 & 2 & 9\\
0 & 1 & 1 & 2\\
2 & 1 & -1 & 0
\end{array}\right]
\]
Usamos el pivote de la columna 1 para anular la entrada del renglón 3.
\[
R_{3}\leftarrow R_{3}-2R_{1}
\]
\[
\left[\begin{array}{ccc|c}
1 & -1 & 2 & 9\\
0 & 1 & 1 & 2\\
0 & 3 & -5 & -18
\end{array}\right]
\]
Usamos el pivote de la columna 2 para anular la entrada del renglón 1.
\[
R_{1}\leftarrow R_{1}+R_{2}
\]
\[
\left[\begin{array}{ccc|c}
1 & 0 & 3 & 11\\
0 & 1 & 1 & 2\\
0 & 3 & -5 & -18
\end{array}\right]
\]
Usamos el pivote de la columna 2 para anular la entrada del renglón 3.
\[
R_{3}\leftarrow R_{3}-3R_{2}
\]
\[
\left[\begin{array}{ccc|c}
1 & 0 & 3 & 11\\
0 & 1 & 1 & 2\\
0 & 0 & -8 & -24
\end{array}\right]
\]
Normalizamos el pivote de la columna 3 para convertirlo en 1.
\[
R_{3}\leftarrow \frac{1}{-8}R_{3}
\]
\[
\left[\begin{array}{ccc|c}
1 & 0 & 3 & 11\\
0 & 1 & 1 & 2\\
0 & 0 & 1 & 3
\end{array}\right]
\]
Usamos el pivote de la columna 3 para anular la entrada del renglón 1.
\[
R_{1}\leftarrow R_{1}-3R_{3}
\]
\[
\left[\begin{array}{ccc|c}
1 & 0 & 0 & 2\\
0 & 1 & 1 & 2\\
0 & 0 & 1 & 3
\end{array}\right]
\]
Usamos el pivote de la columna 3 para anular la entrada del renglón 2.
\[
R_{2}\leftarrow R_{2}-R_{3}
\]
\[
\left[\begin{array}{ccc|c}
1 & 0 & 0 & 2\\
0 & 1 & 0 & -1\\
0 & 0 & 1 & 3
\end{array}\right]
\]
Al finalizar obtenemos la forma escalonada reducida. La solución puede leerse directamente de la última columna.
\[
\boxed{x_{1}=2,\quad x_{2}=-1,\quad x_{3}=3}
\]

\end{Ejem}
%==========================================

\begin{Ejem}
Partimos del sistema
\[
\begin{aligned}
x_{1} + x_{2} - x_{4} &= -3\\
2 x_{1} - x_{2} + x_{3} + x_{4} &= 8\\
x_{1} + 2 x_{2} - x_{3} + 2 x_{4} &= 3\\
3 x_{1} + 2 x_{3} - x_{4} &= 4
\end{aligned}
\]
La matriz aumentada correspondiente es
\[
\left[\begin{array}{cccc|c}
1 & 1 & 0 & -1 & -3\\
2 & -1 & 1 & 1 & 8\\
1 & 2 & -1 & 2 & 3\\
3 & 0 & 2 & -1 & 4
\end{array}\right]
\]
Aplicaremos Gauss--Jordan de manera sistemática. En cada etapa indicaremos qué entrada queremos convertir en pivote y qué entradas deseamos anular.


Usamos el pivote de la columna 1 para anular la entrada del renglón 2.
\[
R_{2}\leftarrow R_{2}-2R_{1}
\]
\[
\left[\begin{array}{cccc|c}
1 & 1 & 0 & -1 & -3\\
0 & -3 & 1 & 3 & 14\\
1 & 2 & -1 & 2 & 3\\
3 & 0 & 2 & -1 & 4
\end{array}\right]
\]
 Usamos el pivote de la columna 1 para anular la entrada del renglón 3.
\[
R_{3}\leftarrow R_{3}-R_{1}
\]
\[
\left[\begin{array}{cccc|c}
1 & 1 & 0 & -1 & -3\\
0 & -3 & 1 & 3 & 14\\
0 & 1 & -1 & 3 & 6\\
3 & 0 & 2 & -1 & 4
\end{array}\right]
\]
 Usamos el pivote de la columna 1 para anular la entrada del renglón 4.
\[
R_{4}\leftarrow R_{4}-3R_{1}
\]
\[
\left[\begin{array}{cccc|c}
1 & 1 & 0 & -1 & -3\\
0 & -3 & 1 & 3 & 14\\
0 & 1 & -1 & 3 & 6\\
0 & -3 & 2 & 2 & 13
\end{array}\right]
\]
Normalizamos el pivote de la columna 2 para convertirlo en 1.
\[
R_{2}\leftarrow \frac{1}{-3}R_{2}
\]
\[
\left[\begin{array}{cccc|c}
1 & 1 & 0 & -1 & -3\\
0 & 1 & - \frac{1}{3} & -1 & - \frac{14}{3}\\
0 & 1 & -1 & 3 & 6\\
0 & -3 & 2 & 2 & 13
\end{array}\right]
\]
Usamos el pivote de la columna 2 para anular la entrada del renglón 1.
\[
R_{1}\leftarrow R_{1}-R_{2}
\]
\[
\left[\begin{array}{cccc|c}
1 & 0 & \frac{1}{3} & 0 & \frac{5}{3}\\
0 & 1 & - \frac{1}{3} & -1 & - \frac{14}{3}\\
0 & 1 & -1 & 3 & 6\\
0 & -3 & 2 & 2 & 13
\end{array}\right]
\]
 Usamos el pivote de la columna 2 para anular la entrada del renglón 3.
\[
R_{3}\leftarrow R_{3}-R_{2}
\]
\[
\left[\begin{array}{cccc|c}
1 & 0 & \frac{1}{3} & 0 & \frac{5}{3}\\
0 & 1 & - \frac{1}{3} & -1 & - \frac{14}{3}\\
0 & 0 & - \frac{2}{3} & 4 & \frac{32}{3}\\
0 & -3 & 2 & 2 & 13
\end{array}\right]
\]
Usamos el pivote de la columna 2 para anular la entrada del renglón 4.
\[
R_{4}\leftarrow R_{4}+3R_{2}
\]
\[
\left[\begin{array}{cccc|c}
1 & 0 & \frac{1}{3} & 0 & \frac{5}{3}\\
0 & 1 & - \frac{1}{3} & -1 & - \frac{14}{3}\\
0 & 0 & - \frac{2}{3} & 4 & \frac{32}{3}\\
0 & 0 & 1 & -1 & -1
\end{array}\right]
\]
Normalizamos el pivote de la columna 3 para convertirlo en 1.
\[
R_{3}\leftarrow \frac{1}{- \frac{2}{3}}R_{3}
\]
\[
\left[\begin{array}{cccc|c}
1 & 0 & \frac{1}{3} & 0 & \frac{5}{3}\\
0 & 1 & - \frac{1}{3} & -1 & - \frac{14}{3}\\
0 & 0 & 1 & -6 & -16\\
0 & 0 & 1 & -1 & -1
\end{array}\right]
\]
Usamos el pivote de la columna 3 para anular la entrada del renglón 1.
\[
R_{1}\leftarrow R_{1}-\frac{1}{3}R_{3}
\]
\[
\left[\begin{array}{cccc|c}
1 & 0 & 0 & 2 & 7\\
0 & 1 & - \frac{1}{3} & -1 & - \frac{14}{3}\\
0 & 0 & 1 & -6 & -16\\
0 & 0 & 1 & -1 & -1
\end{array}\right]
\]
 Usamos el pivote de la columna 3 para anular la entrada del renglón 2.
\[
R_{2}\leftarrow R_{2}+\frac{1}{3}R_{3}
\]
\[
\left[\begin{array}{cccc|c}
1 & 0 & 0 & 2 & 7\\
0 & 1 & 0 & -3 & -10\\
0 & 0 & 1 & -6 & -16\\
0 & 0 & 1 & -1 & -1
\end{array}\right]
\]
 Usamos el pivote de la columna 3 para anular la entrada del renglón 4.
\[
R_{4}\leftarrow R_{4}-R_{3}
\]
\[
\left[\begin{array}{cccc|c}
1 & 0 & 0 & 2 & 7\\
0 & 1 & 0 & -3 & -10\\
0 & 0 & 1 & -6 & -16\\
0 & 0 & 0 & 5 & 15
\end{array}\right]
\]
Normalizamos el pivote de la columna 4 para convertirlo en 1.
\[
R_{4}\leftarrow \frac{1}{5}R_{4}
\]
\[
\left[\begin{array}{cccc|c}
1 & 0 & 0 & 2 & 7\\
0 & 1 & 0 & -3 & -10\\
0 & 0 & 1 & -6 & -16\\
0 & 0 & 0 & 1 & 3
\end{array}\right]
\]
 Usamos el pivote de la columna 4 para anular la entrada del renglón 1.
\[
R_{1}\leftarrow R_{1}-2R_{4}
\]
\[
\left[\begin{array}{cccc|c}
1 & 0 & 0 & 0 & 1\\
0 & 1 & 0 & -3 & -10\\
0 & 0 & 1 & -6 & -16\\
0 & 0 & 0 & 1 & 3
\end{array}\right]
\]
Usamos el pivote de la columna 4 para anular la entrada del renglón 2.
\[
R_{2}\leftarrow R_{2}+3R_{4}
\]
\[
\left[\begin{array}{cccc|c}
1 & 0 & 0 & 0 & 1\\
0 & 1 & 0 & 0 & -1\\
0 & 0 & 1 & -6 & -16\\
0 & 0 & 0 & 1 & 3
\end{array}\right]
\]
Usamos el pivote de la columna 4 para anular la entrada del renglón 3.
\[
R_{3}\leftarrow R_{3}+6R_{4}
\]
\[
\left[\begin{array}{cccc|c}
1 & 0 & 0 & 0 & 1\\
0 & 1 & 0 & 0 & -1\\
0 & 0 & 1 & 0 & 2\\
0 & 0 & 0 & 1 & 3
\end{array}\right]
\]
Al finalizar obtenemos la forma escalonada reducida. La solución puede leerse directamente de la última columna.
\[
\boxed{x_{1}=1,\quad x_{2}=-1,\quad x_{3}=2,\quad x_{4}=3}
\]

\end{Ejem}
%==========================================
\begin{Ejem}
Partimos del sistema
\[
\begin{aligned}
x_{1} + x_{2} - x_{3} + 2 x_{4} + x_{5} &= 8\\
2 x_{1} + 3 x_{2} + x_{3} + x_{4} - 2 x_{5} &= 4\\
- x_{1} + 2 x_{2} + 3 x_{3} - x_{4} + x_{5} &= 1\\
3 x_{1} - x_{2} + 2 x_{3} + 2 x_{4} + x_{5} &= 3\\
2 x_{1} + x_{2} + x_{3} - x_{4} + 3 x_{5} &= 8
\end{aligned}
\]
La matriz aumentada correspondiente es
\[
\left[\begin{array}{ccccc|c}
1 & 1 & -1 & 2 & 1 & 8\\
2 & 3 & 1 & 1 & -2 & 4\\
-1 & 2 & 3 & -1 & 1 & 1\\
3 & -1 & 2 & 2 & 1 & 3\\
2 & 1 & 1 & -1 & 3 & 8
\end{array}\right]
\]
Aplicaremos Gauss--Jordan de manera sistemática. En cada etapa indicaremos qué entrada queremos convertir en pivote y qué entradas deseamos anular. Usamos el pivote de la columna 1 para anular la entrada del renglón 2.
\[
R_{2}\leftarrow R_{2}-2R_{1}
\]
\[
\left[\begin{array}{ccccc|c}
1 & 1 & -1 & 2 & 1 & 8\\
0 & 1 & 3 & -3 & -4 & -12\\
-1 & 2 & 3 & -1 & 1 & 1\\
3 & -1 & 2 & 2 & 1 & 3\\
2 & 1 & 1 & -1 & 3 & 8
\end{array}\right]
\]
Usamos el pivote de la columna 1 para anular la entrada del renglón 3.
\[
R_{3}\leftarrow R_{3}+R_{1}
\]
\[
\left[\begin{array}{ccccc|c}
1 & 1 & -1 & 2 & 1 & 8\\
0 & 1 & 3 & -3 & -4 & -12\\
0 & 3 & 2 & 1 & 2 & 9\\
3 & -1 & 2 & 2 & 1 & 3\\
2 & 1 & 1 & -1 & 3 & 8
\end{array}\right]
\]
Usamos el pivote de la columna 1 para anular la entrada del renglón 4.
\[
R_{4}\leftarrow R_{4}-3R_{1}
\]
\[
\left[\begin{array}{ccccc|c}
1 & 1 & -1 & 2 & 1 & 8\\
0 & 1 & 3 & -3 & -4 & -12\\
0 & 3 & 2 & 1 & 2 & 9\\
0 & -4 & 5 & -4 & -2 & -21\\
2 & 1 & 1 & -1 & 3 & 8
\end{array}\right]
\]
Usamos el pivote de la columna 1 para anular la entrada del renglón 5.
\[
R_{5}\leftarrow R_{5}-2R_{1}
\]
\[
\left[\begin{array}{ccccc|c}
1 & 1 & -1 & 2 & 1 & 8\\
0 & 1 & 3 & -3 & -4 & -12\\
0 & 3 & 2 & 1 & 2 & 9\\
0 & -4 & 5 & -4 & -2 & -21\\
0 & -1 & 3 & -5 & 1 & -8
\end{array}\right]
\]
 Usamos el pivote de la columna 2 para anular la entrada del renglón 1.
\[
R_{1}\leftarrow R_{1}-R_{2}
\]
\[
\left[\begin{array}{ccccc|c}
1 & 0 & -4 & 5 & 5 & 20\\
0 & 1 & 3 & -3 & -4 & -12\\
0 & 3 & 2 & 1 & 2 & 9\\
0 & -4 & 5 & -4 & -2 & -21\\
0 & -1 & 3 & -5 & 1 & -8
\end{array}\right]
\]
Usamos el pivote de la columna 2 para anular la entrada del renglón 3.
\[
R_{3}\leftarrow R_{3}-3R_{2}
\]
\[
\left[\begin{array}{ccccc|c}
1 & 0 & -4 & 5 & 5 & 20\\
0 & 1 & 3 & -3 & -4 & -12\\
0 & 0 & -7 & 10 & 14 & 45\\
0 & -4 & 5 & -4 & -2 & -21\\
0 & -1 & 3 & -5 & 1 & -8
\end{array}\right]
\]

 Usamos el pivote de la columna 2 para anular la entrada del renglón 4.
\[
R_{4}\leftarrow R_{4}+4R_{2}
\]
\[
\left[\begin{array}{ccccc|c}
1 & 0 & -4 & 5 & 5 & 20\\
0 & 1 & 3 & -3 & -4 & -12\\
0 & 0 & -7 & 10 & 14 & 45\\
0 & 0 & 17 & -16 & -18 & -69\\
0 & -1 & 3 & -5 & 1 & -8
\end{array}\right]
\]
Usamos el pivote de la columna 2 para anular la entrada del renglón 5.
\[
R_{5}\leftarrow R_{5}+R_{2}
\]
\[
\left[\begin{array}{ccccc|c}
1 & 0 & -4 & 5 & 5 & 20\\
0 & 1 & 3 & -3 & -4 & -12\\
0 & 0 & -7 & 10 & 14 & 45\\
0 & 0 & 17 & -16 & -18 & -69\\
0 & 0 & 6 & -8 & -3 & -20
\end{array}\right]
\]
Normalizamos el pivote de la columna 3 para convertirlo en 1.
\[
R_{3}\leftarrow \frac{1}{-7}R_{3}
\]
\[
\left[\begin{array}{ccccc|c}
1 & 0 & -4 & 5 & 5 & 20\\
0 & 1 & 3 & -3 & -4 & -12\\
0 & 0 & 1 & - \frac{10}{7} & -2 & - \frac{45}{7}\\
0 & 0 & 17 & -16 & -18 & -69\\
0 & 0 & 6 & -8 & -3 & -20
\end{array}\right]
\]
Usamos el pivote de la columna 3 para anular la entrada del renglón 1.
\[
R_{1}\leftarrow R_{1}+4R_{3}
\]
\[
\left[\begin{array}{ccccc|c}
1 & 0 & 0 & - \frac{5}{7} & -3 & - \frac{40}{7}\\
0 & 1 & 3 & -3 & -4 & -12\\
0 & 0 & 1 & - \frac{10}{7} & -2 & - \frac{45}{7}\\
0 & 0 & 17 & -16 & -18 & -69\\
0 & 0 & 6 & -8 & -3 & -20
\end{array}\right]
\]
 Usamos el pivote de la columna 3 para anular la entrada del renglón 2.
\[
R_{2}\leftarrow R_{2}-3R_{3}
\]
\[
\left[\begin{array}{ccccc|c}
1 & 0 & 0 & - \frac{5}{7} & -3 & - \frac{40}{7}\\
0 & 1 & 0 & \frac{9}{7} & 2 & \frac{51}{7}\\
0 & 0 & 1 & - \frac{10}{7} & -2 & - \frac{45}{7}\\
0 & 0 & 17 & -16 & -18 & -69\\
0 & 0 & 6 & -8 & -3 & -20
\end{array}\right]
\]
Usamos el pivote de la columna 3 para anular la entrada del renglón 4.
\[
R_{4}\leftarrow R_{4}-17R_{3}
\]
\[
\left[\begin{array}{ccccc|c}
1 & 0 & 0 & - \frac{5}{7} & -3 & - \frac{40}{7}\\
0 & 1 & 0 & \frac{9}{7} & 2 & \frac{51}{7}\\
0 & 0 & 1 & - \frac{10}{7} & -2 & - \frac{45}{7}\\
0 & 0 & 0 & \frac{58}{7} & 16 & \frac{282}{7}\\
0 & 0 & 6 & -8 & -3 & -20
\end{array}\right]
\]
 Usamos el pivote de la columna 3 para anular la entrada del renglón 5.
\[
R_{5}\leftarrow R_{5}-6R_{3}
\]
\[
\left[\begin{array}{ccccc|c}
1 & 0 & 0 & - \frac{5}{7} & -3 & - \frac{40}{7}\\
0 & 1 & 0 & \frac{9}{7} & 2 & \frac{51}{7}\\
0 & 0 & 1 & - \frac{10}{7} & -2 & - \frac{45}{7}\\
0 & 0 & 0 & \frac{58}{7} & 16 & \frac{282}{7}\\
0 & 0 & 0 & \frac{4}{7} & 9 & \frac{130}{7}
\end{array}\right]
\]
Normalizamos el pivote de la columna 4 para convertirlo en 1.
\[
R_{4}\leftarrow \frac{1}{\frac{58}{7}}R_{4}
\]
\[
\left[\begin{array}{ccccc|c}
1 & 0 & 0 & - \frac{5}{7} & -3 & - \frac{40}{7}\\
0 & 1 & 0 & \frac{9}{7} & 2 & \frac{51}{7}\\
0 & 0 & 1 & - \frac{10}{7} & -2 & - \frac{45}{7}\\
0 & 0 & 0 & 1 & \frac{56}{29} & \frac{141}{29}\\
0 & 0 & 0 & \frac{4}{7} & 9 & \frac{130}{7}
\end{array}\right]
\]
 Usamos el pivote de la columna 4 para anular la entrada del renglón 1.
\[
R_{1}\leftarrow R_{1}+\frac{5}{7}R_{4}
\]
\[
\left[\begin{array}{ccccc|c}
1 & 0 & 0 & 0 & - \frac{47}{29} & - \frac{65}{29}\\
0 & 1 & 0 & \frac{9}{7} & 2 & \frac{51}{7}\\
0 & 0 & 1 & - \frac{10}{7} & -2 & - \frac{45}{7}\\
0 & 0 & 0 & 1 & \frac{56}{29} & \frac{141}{29}\\
0 & 0 & 0 & \frac{4}{7} & 9 & \frac{130}{7}
\end{array}\right]
\]
Usamos el pivote de la columna 4 para anular la entrada del renglón 2.
\[
R_{2}\leftarrow R_{2}-\frac{9}{7}R_{4}
\]
\[
\left[\begin{array}{ccccc|c}
1 & 0 & 0 & 0 & - \frac{47}{29} & - \frac{65}{29}\\
0 & 1 & 0 & 0 & - \frac{14}{29} & \frac{30}{29}\\
0 & 0 & 1 & - \frac{10}{7} & -2 & - \frac{45}{7}\\
0 & 0 & 0 & 1 & \frac{56}{29} & \frac{141}{29}\\
0 & 0 & 0 & \frac{4}{7} & 9 & \frac{130}{7}
\end{array}\right]
\]
 Usamos el pivote de la columna 4 para anular la entrada del renglón 3.
\[
R_{3}\leftarrow R_{3}+\frac{10}{7}R_{4}
\]
\[
\left[\begin{array}{ccccc|c}
1 & 0 & 0 & 0 & - \frac{47}{29} & - \frac{65}{29}\\
0 & 1 & 0 & 0 & - \frac{14}{29} & \frac{30}{29}\\
0 & 0 & 1 & 0 & \frac{22}{29} & \frac{15}{29}\\
0 & 0 & 0 & 1 & \frac{56}{29} & \frac{141}{29}\\
0 & 0 & 0 & \frac{4}{7} & 9 & \frac{130}{7}
\end{array}\right]
\]
 Usamos el pivote de la columna 4 para anular la entrada del renglón 5.
\[
R_{5}\leftarrow R_{5}-\frac{4}{7}R_{4}
\]
\[
\left[\begin{array}{ccccc|c}
1 & 0 & 0 & 0 & - \frac{47}{29} & - \frac{65}{29}\\
0 & 1 & 0 & 0 & - \frac{14}{29} & \frac{30}{29}\\
0 & 0 & 1 & 0 & \frac{22}{29} & \frac{15}{29}\\
0 & 0 & 0 & 1 & \frac{56}{29} & \frac{141}{29}\\
0 & 0 & 0 & 0 & \frac{229}{29} & \frac{458}{29}
\end{array}\right]
\]
 Normalizamos el pivote de la columna 5 para convertirlo en 1.
\[
R_{5}\leftarrow \frac{1}{\frac{229}{29}}R_{5}
\]
\[
\left[\begin{array}{ccccc|c}
1 & 0 & 0 & 0 & - \frac{47}{29} & - \frac{65}{29}\\
0 & 1 & 0 & 0 & - \frac{14}{29} & \frac{30}{29}\\
0 & 0 & 1 & 0 & \frac{22}{29} & \frac{15}{29}\\
0 & 0 & 0 & 1 & \frac{56}{29} & \frac{141}{29}\\
0 & 0 & 0 & 0 & 1 & 2
\end{array}\right]
\]
Usamos el pivote de la columna 5 para anular la entrada del renglón 1.
\[
R_{1}\leftarrow R_{1}+\frac{47}{29}R_{5}
\]
\[
\left[\begin{array}{ccccc|c}
1 & 0 & 0 & 0 & 0 & 1\\
0 & 1 & 0 & 0 & - \frac{14}{29} & \frac{30}{29}\\
0 & 0 & 1 & 0 & \frac{22}{29} & \frac{15}{29}\\
0 & 0 & 0 & 1 & \frac{56}{29} & \frac{141}{29}\\
0 & 0 & 0 & 0 & 1 & 2
\end{array}\right]
\]
Usamos el pivote de la columna 5 para anular la entrada del renglón 2.
\[
R_{2}\leftarrow R_{2}+\frac{14}{29}R_{5}
\]
\[
\left[\begin{array}{ccccc|c}
1 & 0 & 0 & 0 & 0 & 1\\
0 & 1 & 0 & 0 & 0 & 2\\
0 & 0 & 1 & 0 & \frac{22}{29} & \frac{15}{29}\\
0 & 0 & 0 & 1 & \frac{56}{29} & \frac{141}{29}\\
0 & 0 & 0 & 0 & 1 & 2
\end{array}\right]
\]
 Usamos el pivote de la columna 5 para anular la entrada del renglón 3.
\[
R_{3}\leftarrow R_{3}-\frac{22}{29}R_{5}
\]
\[
\left[\begin{array}{ccccc|c}
1 & 0 & 0 & 0 & 0 & 1\\
0 & 1 & 0 & 0 & 0 & 2\\
0 & 0 & 1 & 0 & 0 & -1\\
0 & 0 & 0 & 1 & \frac{56}{29} & \frac{141}{29}\\
0 & 0 & 0 & 0 & 1 & 2
\end{array}\right]
\]
Usamos el pivote de la columna 5 para anular la entrada del renglón 4.
\[
R_{4}\leftarrow R_{4}-\frac{56}{29}R_{5}
\]
\[
\left[\begin{array}{ccccc|c}
1 & 0 & 0 & 0 & 0 & 1\\
0 & 1 & 0 & 0 & 0 & 2\\
0 & 0 & 1 & 0 & 0 & -1\\
0 & 0 & 0 & 1 & 0 & 1\\
0 & 0 & 0 & 0 & 1 & 2
\end{array}\right]
\]
Al finalizar obtenemos la forma escalonada reducida. La solución puede leerse directamente de la última columna.
\[
\boxed{x_{1}=1,\quad x_{2}=2,\quad x_{3}=-1,\quad x_{4}=1,\quad x_{5}=2}
\]
\end{Ejem}
%==========================================
\begin{Ejem}

Consideremos
\[
\begin{aligned}
x_{1} + 2 x_{2} - x_{3} + x_{4} &= 4\\
2 x_{1} + 4 x_{2} - 2 x_{3} + 2 x_{4} &= 8\\
x_{2} + x_{3} - x_{4} &= 1\\
x_{1} + 3 x_{2} &= 5
\end{aligned}
\]
Su matriz aumentada es
\[
\left[\begin{array}{cccc|c}
1 & 2 & -1 & 1 & 4\\
2 & 4 & -2 & 2 & 8\\
0 & 1 & 1 & -1 & 1\\
1 & 3 & 0 & 0 & 5
\end{array}\right]
\]
y al reducirla por renglones obtenemos
\[
\left[\begin{array}{cccc|c}
1 & 0 & -3 & 3 & 2\\
0 & 1 & 1 & -1 & 1\\
0 & 0 & 0 & 0 & 0\\
0 & 0 & 0 & 0 & 0
\end{array}\right]
\]
Las columnas 1 y 2 contienen pivotes; por tanto, \(x_1\) y \(x_2\) son variables básicas. Las variables
\(x_3\) y \(x_4\) son libres. Tomemos
\[
x_3=s,\qquad x_4=t.
\]
De los dos renglones no nulos obtenemos
\[
x_1-3x_3+3x_4=2,\qquad x_2+x_3-x_4=1.
\]
Así,
\[
x_1=2+3s-3t,\qquad x_2=1-s+t.
\]
Por tanto,
\[
\boxed{(x_1,x_2,x_3,x_4)=(2+3s-3t,\,1-s+t,\,s,\,t),\qquad s,t\in\R.}
\]
La presencia de dos variables libres produce una familia bidimensional de soluciones.
\end{Ejem}
%==========================================
\begin{Ejem}

Identificar la contradicción algebraica que certifica que un sistema no tiene solución. Consideremos
\[
\begin{aligned}
x_{1} + x_{2} - x_{3} &= 2\\
2 x_{1} + 2 x_{2} - 2 x_{3} &= 5\\
x_{1} - x_{2} + x_{3} &= 0
\end{aligned}
\]
Al aplicar operaciones elementales se obtiene una forma reducida que contiene un renglón equivalente a
\[
0x_1+0x_2+0x_3=1.
\]
Es decir,
\[
0=1,
\]
lo cual es imposible. Por tanto,
\[
\boxed{\text{el sistema es incompatible y no tiene solución}.}
\]
Lo importante es que la incompatibilidad se reconoce durante la reducción; no es necesario continuar buscando valores para las incógnitas.
\end{Ejem}

\begin{Nota}

Gauss busca una matriz escalonada y termina con sustitución hacia atrás. Gauss--Jordan continúa hasta la forma escalonada reducida. Para resolver manualmente un sistema grande, Gauss suele requerir menos operaciones. Gauss--Jordan, en cambio, proporciona una descripción muy transparente de variables básicas, variables libres, rango y consistencia.
\end{Nota}
%<<>>==<<>><<>>==<<>><<>>==<<>><<>>==<<>><<>>==<<>><<>>
\section{Sistemas con parámetros}
%<<>>==<<>><<>>==<<>><<>>==<<>><<>>==<<>><<>>==<<>><<>>
Los parámetros permiten estudiar familias completas de sistemas.

%<<>>==<<>><<>>==<<>><<>>==<<>><<>>==<<>><<>>==<<>><<>>
\subsection*{Parámetros como variables libres}
%<<>>==<<>><<>>==<<>><<>>==<<>><<>>==<<>><<>>==<<>><<>>
Consideremos
\[
\begin{cases}
x+y+z=4,\\
2x+2y+2z=8.
\end{cases}
\]
Después de reducir,
\[
x+y+z=4.
\]
Tomando
\[
y=s,\qquad z=t,
\]
obtenemos
\[
x=4-s-t.
\]
Aquí $s$ y $t$ son parámetros que describen todas las soluciones.
%<<>>==<<>><<>>==<<>><<>>==<<>><<>>==<<>><<>>==<<>><<>>
\subsection*{Parámetros en los coeficientes}
%<<>>==<<>><<>>==<<>><<>>==<<>><<>>==<<>><<>>==<<>><<>>
Consideremos
\[
\begin{cases}
x+y=2,\\
2x+ay=4.
\end{cases}
\]
La matriz aumentada es
\[
\left[
\begin{array}{cc|c}
1&1&2\\
2&a&4
\end{array}
\right].
\]
Aplicando
\[
R_2\leftarrow R_2-2R_1
\]
obtenemos
\[
\left[
\begin{array}{cc|c}
1&1&2\\
0&a-2&0
\end{array}
\right].
\]

Si $a\neq2$, entonces
\[
(a-2)y=0\quad\Longrightarrow\quad y=0,
\]
y $x=2$. Existe una única solución.

Si $a=2$, el segundo renglón se vuelve nulo:
\[
[0\quad0\mid0],
\]
y queda una ecuación con dos incógnitas. Existen infinitas soluciones. Este ejemplo muestra por qué, al trabajar con parámetros, no debemos dividir por una expresión como $a-2$ sin estudiar antes el caso en que puede ser cero.
%<<>>==<<>><<>>==<<>><<>>==<<>><<>>==<<>><<>>==<<>><<>>
\subsection*{Un parámetro que produce incompatibilidad}
%<<>>==<<>><<>>==<<>><<>>==<<>><<>>==<<>><<>>==<<>><<>>
\[
\begin{cases}
x+y=2,\\
2x+2y=k.
\end{cases}
\]
Reduciendo:
\[
\left[
\begin{array}{cc|c}
1&1&2\\
0&0&k-4
\end{array}
\right].
\]
Por tanto:
\[
\begin{cases}
k=4 &\Longrightarrow \text{infinitas soluciones},\\
k\neq4 &\Longrightarrow \text{ninguna solución}.
\end{cases}
\]
%<<>>==<<>><<>>==<<>><<>>==<<>><<>>==<<>><<>>==<<>><<>>
\section{Matrices: Introducci\'on}
%<<>>==<<>><<>>==<<>><<>>==<<>><<>>==<<>><<>>==<<>><<>>
\begin{definicion}
El \textbf{rango} de una matriz $A$, denotado por $\rank(A)$, es el número de pivotes de una forma escalonada de $A$.
\end{definicion}

Equivalentemente, es el número de renglones no nulos de su forma escalonada.

\begin{ejemplo}
Si
\[
A\sim
\begin{pmatrix}
1&2&0&3\\
0&1&4&2\\
0&0&0&0
\end{pmatrix},
\]
entonces
\[
\rank(A)=2.
\]
\end{ejemplo}
%<<>>==<<>><<>>==<<>><<>>==<<>><<>>==<<>><<>>==<<>><<>>
\subsection*{Rango y consistencia}
%<<>>==<<>><<>>==<<>><<>>==<<>><<>>==<<>><<>>==<<>><<>>
Para el sistema
\[
A\mathbf{x}=\mathbf b,
\]
debemos comparar la matriz de coeficientes $A$ con la matriz aumentada $[A\mid\mathbf b]$.

\begin{teorema}[Criterio de consistencia]
El sistema $A\mathbf{x}=\mathbf b$ es compatible si y sólo si
\[
\rank(A)=\rank([A\mid\mathbf b]).
\]
\end{teorema}

\begin{observacion}
Si al reducir la matriz aumentada aparece un renglón
\[
[0\ \cdots\ 0\mid c],\qquad c\neq0,
\]
entonces la última columna introduce un pivote que no existe en $A$. Por ello
\[
\rank(A)<\rank([A\mid\mathbf b]),
\]
y el sistema es incompatible.
\end{observacion}

Si el sistema es compatible y tiene $n$ incógnitas:
\[
\begin{cases}
\rank(A)=n &\Longrightarrow \text{solución única},\\
\rank(A)<n &\Longrightarrow \text{infinitas soluciones}.
\end{cases}
\]
%<<>>==<<>><<>>==<<>><<>>==<<>><<>>==<<>><<>>==<<>><<>>
\subsection*{Solución mediante matriz inversa}
%<<>>==<<>><<>>==<<>><<>>==<<>><<>>==<<>><<>>==<<>><<>>
Supongamos que $A$ es una matriz cuadrada e invertible. Si
\[
A\mathbf{x}=\mathbf b,
\]
multiplicamos por $A^{-1}$:
\[
A^{-1}A\mathbf{x}=A^{-1}\mathbf b.
\]
Como
\[
A^{-1}A=I,
\]
obtenemos
\[
\boxed{\mathbf{x}=A^{-1}\mathbf b.}
\]

\begin{ejemplo}
Sea
\[
A=
\begin{pmatrix}
2&1\\
1&1
\end{pmatrix},
\qquad
\mathbf b=
\begin{pmatrix}5\\3\end{pmatrix}.
\]
Como
\[
A^{-1}=
\begin{pmatrix}
1&-1\\
-1&2
\end{pmatrix},
\]
se tiene
\[
\mathbf x=A^{-1}\mathbf b
=
\begin{pmatrix}
1&-1\\
-1&2
\end{pmatrix}
\begin{pmatrix}5\\3\end{pmatrix}
=
\begin{pmatrix}2\\1\end{pmatrix}.
\]
\end{ejemplo}

\begin{observacion}
La fórmula $A^{-1}\mathbf b$ sólo tiene sentido cuando $A$ es cuadrada e invertible. Además, calcular explícitamente la inversa no suele ser el procedimiento más eficiente para resolver un único sistema. El método es conceptualmente importante porque relaciona invertibilidad con existencia y unicidad.
\end{observacion}
%<<>>==<<>><<>>==<<>><<>>==<<>><<>>==<<>><<>>==<<>><<>>
\subsection*{Regla de Cramer}
%<<>>==<<>><<>>==<<>><<>>==<<>><<>>==<<>><<>>==<<>><<>>
La regla de Cramer proporciona una fórmula para sistemas cuadrados cuya matriz de coeficientes tiene determinante distinto de cero.

Sea
\[
A\mathbf x=\mathbf b,\qquad \det(A)\neq0.
\]
Para cada $i$, sea $A_i$ la matriz obtenida al sustituir la columna $i$ de $A$ por $\mathbf b$. Entonces
\[
\boxed{x_i=\frac{\det(A_i)}{\det(A)}}.
\]

\begin{ejemplo}
Resolvamos
\[
\begin{cases}
2x+y=5,\\
x-y=1.
\end{cases}
\]
Tenemos
\[
A=
\begin{pmatrix}2&1\\1&-1\end{pmatrix},
\qquad
\det(A)=-3.
\]
Para $x$,
\[
A_1=
\begin{pmatrix}5&1\\1&-1\end{pmatrix},
\qquad
\det(A_1)=-6,
\]
por lo que
\[
x=\frac{-6}{-3}=2.
\]
Para $y$,
\[
A_2=
\begin{pmatrix}2&5\\1&1\end{pmatrix},
\qquad
\det(A_2)=-3,
\]
y
\[
y=\frac{-3}{-3}=1.
\]
\end{ejemplo}

\begin{observacion}
Cramer requiere $\det(A)\neq0$. Es útil teóricamente y para sistemas pequeños, pero para sistemas grandes la eliminación gaussiana suele ser más eficiente.
\end{observacion}
%<<>>==<<>><<>>==<<>><<>>==<<>><<>>==<<>><<>>==<<>><<>>
\section{Aplicaciones}
%<<>>==<<>><<>>==<<>><<>>==<<>><<>>==<<>><<>>==<<>><<>>
%<<>>==<<>><<>>==<<>><<>>==<<>><<>>==<<>><<>>==<<>><<>>
\subsection*{Mezclas}
%<<>>==<<>><<>>==<<>><<>>==<<>><<>>==<<>><<>>==<<>><<>>
Una solución de $10$ litros debe contener $30\%$ de cierta sustancia. Se dispone de una solución al
$20\%$ y otra al $50\%$. Si $x$ y $y$ son los litros utilizados de cada una,
\[
x+y=10.
\]
La cantidad de sustancia pura debe satisfacer
\[
0.20x+0.50y=0.30(10)=3.
\]
Así,
\[
\begin{cases}
x+y=10,\\
0.2x+0.5y=3.
\end{cases}
\]
Multiplicando la segunda ecuación por $10$:
\[
2x+5y=30.
\]
De $x=10-y$:
\[
2(10-y)+5y=30,
\]
por lo que $3y=10$ y
\[
y=\frac{10}{3},\qquad x=\frac{20}{3}.
\]
%<<>>==<<>><<>>==<<>><<>>==<<>><<>>==<<>><<>>==<<>><<>>
\subsection*{Circuitos}
%<<>>==<<>><<>>==<<>><<>>==<<>><<>>==<<>><<>>==<<>><<>>
En muchos circuitos, las leyes de Kirchhoff producen sistemas lineales en las corrientes desconocidas.
Por ejemplo,
\[
\begin{cases}
I_1-I_2-I_3=0,\\
4I_1+2I_2=12,\\
2I_2-3I_3=0.
\end{cases}
\]
La primera ecuación puede representar conservación de corriente en un nodo; las restantes, balances de
voltaje en mallas. Una vez formulado el modelo, el problema se convierte en resolver un sistema lineal.
%<<>>==<<>><<>>==<<>><<>>==<<>><<>>==<<>><<>>==<<>><<>>
\subsection*{Equilibrio de fuerzas}
%<<>>==<<>><<>>==<<>><<>>==<<>><<>>==<<>><<>>==<<>><<>>
Si varias fuerzas actúan sobre una estructura en equilibrio, la suma de componentes debe ser cero:
\[
\sum F_x=0,\qquad \sum F_y=0.
\]
Cuando las magnitudes desconocidas aparecen linealmente, estas ecuaciones forman un sistema lineal.
%<<>>==<<>><<>>==<<>><<>>==<<>><<>>==<<>><<>>==<<>><<>>
\section{Ejemplos}
%<<>>==<<>><<>>==<<>><<>>==<<>><<>>==<<>><<>>==<<>><<>>
%<<>>==<<>><<>>==<<>><<>>==<<>><<>>==<<>><<>>==<<>><<>>
%<<>>==<<>><<>>==<<>><<>>==<<>><<>>==<<>><<>>==<<>><<>>
\begin{Ejem} Resolver el siguiente sistema de ecuaciones lineales
\[
\begin{cases}
x+2y-z=1,\\
2x+4y-2z=2,\\
x+y+z=3.
\end{cases}
\]
La matriz aumentada es
\[
\left[
\begin{array}{ccc|c}
1&2&-1&1\\
2&4&-2&2\\
1&1&1&3
\end{array}
\right].
\]
Aplicamos
\[
R_2\leftarrow R_2-2R_1,\qquad
R_3\leftarrow R_3-R_1,
\]
y obtenemos
\[
\left[
\begin{array}{ccc|c}
1&2&-1&1\\
0&0&0&0\\
0&-1&2&2
\end{array}
\right].
\]
Intercambiando los dos últimos renglones:
\[
\left[
\begin{array}{ccc|c}
1&2&-1&1\\
0&-1&2&2\\
0&0&0&0
\end{array}
\right].
\]
Hay dos pivotes y tres incógnitas. Por tanto existe una variable libre. Tomemos $z=t$. La segunda
ecuación da
\[
-y+2t=2\quad\Longrightarrow\quad y=2t-2.
\]
La primera:
\[
x+2(2t-2)-t=1,
\]
de donde
\[
x=5-3t.
\]
Así,
\[
\boxed{(x,y,z)=(5-3t,\,2t-2,\,t),\quad t\in\R.}
\]
El sistema es compatible indeterminado.
\end{Ejem}
%<<>>==<<>><<>>==<<>><<>>==<<>><<>>==<<>><<>>==<<>><<>>
%<<>>==<<>><<>>==<<>><<>>==<<>><<>>==<<>><<>>==<<>><<>>
\begin{Ejem} Resolver el siguiente sistema de ecuaciones lineales
\[
\begin{cases}
x+y+z=1,\\
2x+2y+2z=2,\\
x+y+z=k.
\end{cases}
\]
Restando la primera ecuación de la tercera obtenemos
\[
0=k-1.
\]
Por tanto:
\[
k=1\Longrightarrow\text{infinitas soluciones},
\]
mientras que
\[
k\neq1\Longrightarrow\text{sistema incompatible}.
\]
\end{Ejem}
%<<>>==<<>><<>>==<<>><<>>==<<>><<>>==<<>><<>>==<<>><<>>
\section{Ejercicios}
%<<>>==<<>><<>>==<<>><<>>==<<>><<>>==<<>><<>>==<<>><<>>
\textbf{Instrucciones} En cada ejercicio: (1) escriba la matriz aumentada; (2) identifique
el pivote actual; (3) registre cada operación elemental; (4) no avance de columna hasta haber decidido qué entradas deben anularse; (5) continúe hasta RREF; (6) interprete el resultado y compruébelo cuando haya solución única.

\begin{Ejer}

\begin{enumerate}
\begin{multicols}{3}
\item \[
\begin{aligned}
3 x_{1} - 2 x_{2} &= 7\\
2 x_{1} - x_{2} &= 5
\end{aligned}
\]
\item 
\[
\begin{aligned}
- 3 x_{1} + 4 x_{2} &= 10\\
- 4 x_{1} &= -8
\end{aligned}
\]
\item
\[
\begin{aligned}
3 x_{1} + 4 x_{2} &= 25\\
3 x_{1} - 2 x_{2} &= 1
\end{aligned}
\]
\item \[
\begin{aligned}
3 x_{1} + 2 x_{2} &= -11\\
- 2 x_{1} - 2 x_{2} &= 8
\end{aligned}
\]
\item \[
\begin{aligned}
3 x_{1} - 2 x_{2} &= -1\\
- 4 x_{1} - 2 x_{2} &= -8
\end{aligned}
\]
\item \[
\begin{aligned}
- x_{2} + 2 x_{3} &= 0\\
3 x_{1} - x_{2} - 4 x_{3} &= 0\\
x_{1} - 4 x_{2} + 2 x_{3} &= -8
\end{aligned}
\]
\end{multicols}
\end{enumerate}

\end{Ejer}
%<<>>==<<>><<>>==<<>><<>>==<<>><<>>==<<>><<>>==<<>><<>>==<<>>
\begin{Ejer} Resolver los siguientes sistemas de ecuaciones lineales
\begin{enumerate}[resume]
\begin{multicols}{3}
\item \[
\begin{aligned}
2 x_{1} + 3 x_{2} + 3 x_{3} &= 8\\
- x_{1} - 4 x_{2} - 2 x_{3} &= -5\\
- 2 x_{1} - 2 x_{2} - 2 x_{3} &= -6
\end{aligned}
\]
\item \[
\begin{aligned}
- 3 x_{1} - 4 x_{2} - 2 x_{3} &= 6\\
x_{2} - x_{3} &= -6\\
- x_{1} - 2 x_{2} - 2 x_{3} &= -2
\end{aligned}
\]
\item \[
\begin{aligned}
- 2 x_{1} - 4 x_{2} - 3 x_{3} &= -17\\
2 x_{1} - 2 x_{2} &= -4\\
4 x_{1} + 3 x_{2} + 4 x_{3} &= 17
\end{aligned}
\]
\item \[
\begin{aligned}
- 4 x_{1} + 4 x_{3} &= 4\\
- x_{1} - 3 x_{2} + 3 x_{3} &= 8\\
2 x_{3} &= -2
\end{aligned}
\]
\item \[
\begin{aligned}
- 2 x_{1} + 3 x_{2} - 3 x_{3} &= -11\\
2 x_{1} - 3 x_{2} - 3 x_{3} &= -1\\
- 4 x_{1} - 4 x_{2} - x_{3} &= -2
\end{aligned}
\]
\item \[
\begin{aligned}
- 3 x_{1} + x_{2} &= -8\\
3 x_{1} - 2 x_{2} + 2 x_{3} &= 16\\
2 x_{2} - x_{3} &= -7
\end{aligned}
\]
\item \[
\begin{aligned}
x_{1} - 2 x_{2} - 4 x_{3} &= 2\\
- 4 x_{1} + 2 x_{2} + x_{3} &= 1\\
- 2 x_{1} - x_{2} + 4 x_{3} &= -5
\end{aligned}
\]

\end{multicols}
\end{enumerate}
\end{Ejer}
%<<>>==<<>><<>>==<<>><<>>==<<>><<>>==<<>><<>>==<<>>
\begin{Ejer} Resolver los siguientes sistemas de ecuaciones lineales
\begin{enumerate}[resume]
\begin{multicols}{2}
\item \[
\begin{aligned}
- x_{1} + 3 x_{2} + 2 x_{3} - x_{4} &= 5\\
x_{1} + x_{2} - 2 x_{3} + 2 x_{4} &= -2\\
- 3 x_{1} - 2 x_{2} - 2 x_{3} + 2 x_{4} &= -1\\
- x_{1} - 2 x_{2} + 3 x_{3} - 3 x_{4} &= 2
\end{aligned}
\]
\item \[
\begin{aligned}
- 3 x_{1} - 2 x_{2} + 2 x_{3} &= 0\\
- 2 x_{1} + 3 x_{2} - 3 x_{3} + 3 x_{4} &= 25\\
- 3 x_{1} + x_{3} - 2 x_{4} &= -5\\
2 x_{1} - 2 x_{2} - 3 x_{3} &= 5
\end{aligned}
\]
\item \[
\begin{aligned}
x_{1} - x_{2} - 3 x_{4} &= -5\\
- 3 x_{1} - 3 x_{2} + 3 x_{3} + 2 x_{4} &= 7\\
- 3 x_{2} - 3 x_{3} - x_{4} &= -8\\
- 2 x_{1} + 3 x_{2} - 3 x_{3} + x_{4} &= -6
\end{aligned}
\]
\item \[
\begin{aligned}
- 3 x_{1} - 2 x_{2} - 2 x_{3} &= 9\\
x_{1} + x_{2} + 3 x_{3} + x_{4} &= -9\\
3 x_{3} + x_{4} &= -9\\
x_{1} - 2 x_{3} + x_{4} &= 3
\end{aligned}
\]
\item \[
\begin{aligned}
2 x_{1} - 2 x_{2} - x_{3} &= -10\\
2 x_{1} - x_{2} - x_{3} + x_{4} &= -7\\
x_{1} - 3 x_{3} - 2 x_{4} &= -4\\
3 x_{2} - 3 x_{3} + 3 x_{4} &= 3
\end{aligned}
\]
\item \[
\begin{aligned}
- x_{1} + 2 x_{2} - x_{3} - 2 x_{4} &= -14\\
3 x_{1} + 2 x_{3} - x_{4} &= 13\\
3 x_{1} - 3 x_{2} - x_{3} - 2 x_{4} &= 8\\
- 3 x_{1} + x_{2} - x_{4} &= -13
\end{aligned}
\]
\item \[
\begin{aligned}
3 x_{2} + 3 x_{3} - 2 x_{4} &= 9\\
2 x_{1} - x_{2} - 2 x_{3} + x_{4} &= -3\\
3 x_{1} + 2 x_{2} + 3 x_{4} &= -7\\
x_{2} - x_{3} &= -5
\end{aligned}
\]

\end{multicols}
\end{enumerate}
\end{Ejer}
%<<>>==<<>><<>>==<<>><<>>==<<>><<>>==<<>><<>>==<<>>
%<<>>==<<>><<>>==<<>><<>>==<<>><<>>==<<>><<>>==<<>>
\begin{Ejer} Resolver los siguientes sistemas de ecuaciones lineales
\begin{enumerate}[resume]
\begin{multicols}{2}
\item 
\[
\begin{aligned}
x_{1} - x_{2} + 2 x_{3} - 2 x_{5} &= 6\\
- 3 x_{1} + 3 x_{2} - x_{3} + x_{4} &= 14\\
- x_{1} - x_{2} + x_{3} - 3 x_{4} + 3 x_{5} &= -3\\
2 x_{1} - 3 x_{2} - 3 x_{3} - 3 x_{4} &= -21\\
3 x_{1} + x_{2} - 2 x_{3} + 2 x_{4} + x_{5} &= -17
\end{aligned}
\]
\item \[
\begin{aligned}
3 x_{1} - 3 x_{2} + x_{3} - 2 x_{4} - x_{5} &= 16\\
- x_{1} - 3 x_{2} + 2 x_{3} + 3 x_{4} - 2 x_{5} &= 1\\
- 2 x_{1} - 2 x_{2} + 3 x_{3} + 3 x_{4} + 2 x_{5} &= -4\\
- 2 x_{1} + x_{2} + 2 x_{3} + x_{4} - 3 x_{5} &= 0\\
2 x_{1} + x_{2} + 2 x_{3} + x_{4} - 3 x_{5} &= 8
\end{aligned}
\]
\item \[
\begin{aligned}
- 3 x_{1} - x_{2} - 2 x_{3} + 2 x_{4} + x_{5} &= 2\\
3 x_{1} + 3 x_{2} + 3 x_{4} - 2 x_{5} &= 1\\
- x_{2} - 3 x_{3} + 3 x_{4} - 2 x_{5} &= -11\\
- 3 x_{1} - 2 x_{2} + 2 x_{3} - x_{4} + 3 x_{5} &= 17\\
- 3 x_{2} + 3 x_{3} - 3 x_{5} &= 0
\end{aligned}
\]
\item \[
\begin{aligned}
3 x_{1} + 3 x_{3} + 3 x_{4} - 3 x_{5} &= 30\\
x_{1} + 3 x_{2} + x_{3} + 2 x_{4} + 3 x_{5} &= 18\\
2 x_{1} + 2 x_{2} + 3 x_{3} + 3 x_{5} &= 18\\
- 3 x_{1} + 3 x_{2} - 3 x_{3} + 2 x_{4} + 3 x_{5} &= -6\\
- 3 x_{1} - 2 x_{2} - 3 x_{3} + 2 x_{4} + 2 x_{5} &= -20
\end{aligned}
\]
\item \[
\begin{aligned}
x_{1} - 3 x_{2} - 2 x_{3} + 2 x_{4} - x_{5} &= -5\\
3 x_{1} + x_{2} + 3 x_{3} + 3 x_{4} - 2 x_{5} &= -14\\
- x_{1} - 2 x_{2} - 2 x_{3} + x_{4} + x_{5} &= 7\\
- 2 x_{1} + x_{2} + 3 x_{3} + 2 x_{4} + x_{5} &= 9\\
3 x_{1} - 2 x_{2} + 2 x_{3} - 2 x_{4} - 3 x_{5} &= -24
\end{aligned}
\]
\item \[
\begin{aligned}
- 2 x_{3} - x_{4} - 3 x_{5} &= 11\\
3 x_{2} + 3 x_{3} + x_{4} - x_{5} &= 4\\
x_{1} - 3 x_{2} - x_{3} - 3 x_{4} + x_{5} &= -17\\
2 x_{5} &= -6\\
3 x_{1} - 2 x_{3} - 3 x_{4} - 2 x_{5} &= 3
\end{aligned}
\]
\item \[
\begin{aligned}
- x_{1} - x_{2} + x_{3} + x_{4} + x_{5} &= 6\\
- x_{2} + x_{3} - x_{4} &= -4\\
2 x_{2} - 2 x_{3} + 2 x_{4} - x_{5} &= 4\\
3 x_{1} + 2 x_{2} + x_{3} - 2 x_{5} &= -12\\
x_{1} - 2 x_{2} + x_{3} - x_{4} - x_{5} &= -7
\end{aligned}
\]
\item \[
\begin{aligned}
- x_{1} - x_{2} - x_{3} + x_{4} - 2 x_{5} &= -6\\
- 2 x_{2} + 2 x_{3} - 3 x_{4} &= 2\\
3 x_{1} - x_{2} + x_{4} - 2 x_{5} &= 1\\
- 2 x_{1} + 2 x_{2} - 3 x_{3} - 2 x_{5} &= -1\\
- 2 x_{1} - x_{2} + 3 x_{3} + 2 x_{4} - 2 x_{5} &= -14
\end{aligned}
\]

\end{multicols}
\end{enumerate}
\end{Ejer}
%<<>>==<<>><<>>==<<>><<>>==<<>><<>>==<<>><<>>==<<>>
%<<>>==<<>><<>>==<<>><<>>==<<>><<>>==<<>><<>>==<<>>
\begin{Ejer} Resolver los siguientes sistemas de ecuaciones lineales
\begin{enumerate}[resume]
\begin{multicols}{2}
\item 
\[
\begin{aligned}
- x_{2} - 3 x_{3} &= 4\\
- 3 x_{1} + x_{2} + 3 x_{3} &= -16\\
- 3 x_{1} + x_{2} + 3 x_{3} &= -16
\end{aligned}
\]
\item \[
\begin{aligned}
2 x_{1} + 3 x_{2} &= 20\\
x_{1} - 2 x_{2} - 3 x_{3} &= -4\\
x_{1} - 2 x_{2} - 3 x_{3} &= -4
\end{aligned}
\]
\item \[
\begin{aligned}
- 3 x_{2} + 2 x_{3} - 2 x_{4} &= -17\\
3 x_{1} + 3 x_{2} - x_{3} - 3 x_{4} &= 0\\
2 x_{1} - 2 x_{4} &= -4\\
2 x_{1} - 2 x_{4} &= -4
\end{aligned}
\]
\item \[
\begin{aligned}
x_{1} + 3 x_{2} - 3 x_{3} + 2 x_{4} &= -13\\
- 2 x_{2} + 2 x_{4} &= -4\\
- 2 x_{1} + 2 x_{2} - x_{3} + 3 x_{4} &= -17\\
- 2 x_{1} + 2 x_{2} - x_{3} + 3 x_{4} &= -17
\end{aligned}
\]
\item \[
\begin{aligned}
2 x_{1} + 2 x_{2} - 3 x_{3} + x_{4} + x_{5} &= 9\\
- 3 x_{1} - 3 x_{2} + x_{3} + 3 x_{4} + 2 x_{5} &= 8\\
- x_{1} - x_{2} - x_{4} - 3 x_{5} &= 3\\
- x_{1} - x_{2} - 2 x_{3} - 3 x_{4} &= -5\\
- x_{1} - x_{2} - 2 x_{3} - 3 x_{4} &= -5
\end{aligned}
\]
\item \[
\begin{aligned}
- 3 x_{1} - 2 x_{2} + 2 x_{3} - x_{4} + 2 x_{5} &= 7\\
2 x_{1} + x_{2} - 3 x_{3} - 2 x_{5} &= -3\\
- x_{1} + 2 x_{2} + 2 x_{4} - 2 x_{5} &= -10\\
2 x_{1} + x_{2} &= -2\\
2 x_{1} + x_{2} &= -2
\end{aligned}
\]
\item \[
\begin{aligned}
- 2 x_{1} + 2 x_{2} &= 8\\
3 x_{1} + x_{2} + 3 x_{3} &= 8\\
3 x_{1} + x_{2} + 3 x_{3} &= 9
\end{aligned}
\]
\item \[
\begin{aligned}
x_{2} - 3 x_{3} &= 2\\
3 x_{1} - 2 x_{2} - 3 x_{3} &= -10\\
3 x_{1} - 2 x_{2} - 3 x_{3} &= -9
\end{aligned}
\]
\item \[
\begin{aligned}
2 x_{2} + 3 x_{3} + 2 x_{4} &= -1\\
- 2 x_{1} + 3 x_{2} - 2 x_{3} &= -12\\
- 3 x_{2} - x_{3} - 2 x_{4} &= 5\\
- 3 x_{2} - x_{3} - 2 x_{4} &= 6
\end{aligned}
\]
\item \[
\begin{aligned}
3 x_{1} + x_{2} - x_{3} + x_{4} &= -2\\
x_{1} + 3 x_{2} + 2 x_{3} &= 4\\
- 2 x_{1} - 2 x_{2} + 2 x_{3} &= -10\\
- 2 x_{1} - 2 x_{2} + 2 x_{3} &= -9
\end{aligned}
\]
\item \[
\begin{aligned}
- 2 x_{1} + 3 x_{2} - 2 x_{3} + 2 x_{4} + x_{5} &= 6\\
- x_{1} - x_{2} + x_{3} + 2 x_{4} - 2 x_{5} &= 5\\
- x_{1} - 2 x_{2} + 3 x_{3} &= 10\\
- 3 x_{1} + 2 x_{2} - 2 x_{3} + x_{4} - 3 x_{5} &= -11\\
- 3 x_{1} + 2 x_{2} - 2 x_{3} + x_{4} - 3 x_{5} &= -10
\end{aligned}
\]
\item \[
\begin{aligned}
x_{1} - x_{2} - 3 x_{4} &= -7\\
- 2 x_{1} + 2 x_{2} + 3 x_{3} - 3 x_{4} &= 5\\
- x_{1} - 3 x_{3} - 2 x_{4} - 2 x_{5} &= -20\\
- 2 x_{1} - 3 x_{2} - 3 x_{3} + x_{4} - 2 x_{5} &= -21\\
- 2 x_{1} - 3 x_{2} - 3 x_{3} + x_{4} - 2 x_{5} &= -20
\end{aligned}
\]
\end{multicols}
\end{enumerate}
\end{Ejer}
%<<>>==<<>><<>>==<<>><<>>==<<>><<>>==<<>><<>>==<<>>

\begin{Ejer}
\begin{enumerate}
\item Determine cuáles de las siguientes ecuaciones son lineales:
\[
3x-2y+z=4,\qquad x^2+y=1,\qquad xy-z=0,\qquad 2x_1-5x_2=7.
\]
\item Explique con sus palabras la diferencia entre un sistema compatible determinado, compatible
indeterminado e incompatible.
\item ¿Por qué todo sistema homogéneo es compatible?
\item ¿Qué información proporciona un pivote?
\item Explique la diferencia entre forma escalonada y forma escalonada reducida.
\end{enumerate}
\end{Ejer}
%<<>>==<<>><<>>==<<>><<>>==<<>><<>>==<<>><<>>==<<>>
\begin{Ejer} Eliminación de Gauss
Resuelva y clasifique:
\begin{enumerate}
\begin{multicols}{2}
\item
\[
\begin{cases}
x+y=7,\\
2x-y=2.
\end{cases}
\]
\item
\[
\begin{cases}
x+y+z=6,\\
2x-y+3z=9,\\
3x+2y-z=4.
\end{cases}
\]
\item
\[
\begin{cases}
x+2y-z=3,\\
2x+4y-2z=6.
\end{cases}
\]
\item
\[
\begin{cases}
x-y=1,\\
2x-2y=5.
\end{cases}
\]
\end{multicols}
\end{enumerate}
\end{Ejer}
%<<>>==<<>><<>>==<<>><<>>==<<>><<>>==<<>><<>>==<<>>
\begin{Ejer}  Gauss--Jordan y variables libres
\begin{enumerate}

\item Reduzca a forma escalonada reducida:
\[
\left[
\begin{array}{ccc|c}
1&2&-1&3\\
2&5&1&8\\
-1&-1&2&-1
\end{array}
\right].
\]
\item Para
\[
\left[
\begin{array}{cccc|c}
1&2&0&-1&4\\
0&0&1&3&2\\
0&0&0&0&0
\end{array}
\right],
\]
identifique variables básicas y libres y describa todas las soluciones.
\end{enumerate}
\end{Ejer}
%<<>>==<<>><<>>==<<>><<>>==<<>><<>>==<<>><<>>==<<>>
\begin{Ejer} 
Clasifique los sistemas según el valor del parámetro:
\begin{enumerate}

\item
\[
\begin{cases}
x+y=1,\\
2x+ay=2.
\end{cases}
\]
\item
\[
\begin{cases}
x-y=3,\\
2x-2y=k.
\end{cases}
\]
\end{enumerate}
\end{Ejer}
%<<>>==<<>><<>>==<<>><<>>==<<>><<>>==<<>><<>>==<<>>
\begin{Ejer} Aplicaciones
\begin{enumerate}
\item Una mezcla de $20$ kg debe contener $35\%$ de un componente. ¿Cuántos kilogramos de mezclas
al $20\%$ y al $50\%$ deben combinarse?
\item Formule y resuelva un sistema para determinar dos números cuya suma es $18$ y cuya diferencia
es $4$.
\item Tres corrientes satisfacen
\[
I_1-I_2-I_3=0,\qquad 5I_1+2I_2=15,\qquad 2I_2-4I_3=0.
\]
Determine $I_1,I_2,I_3$.
\end{enumerate}
\end{Ejer}
%<<>>==<<>><<>>==<<>><<>>==<<>><<>>==<<>><<>>==<<>>
\begin{Ejer}

Utilice el método de eliminación de Gauss--Jordan para encontrar, si existen, todas las soluciones de los sistemas dados.

\begin{enumerate}[resume]
\begin{multicols}{2}
\item \sistema{
9x_1+9x_2-7x_3&=6,\\
-7x_1-x_3&=-10,\\
9x_1+6x_2+8x_3&=45.
}

\item \sistema{
x_1-2x_2+3x_3&=11,\\
4x_1+x_2-x_3&=4,\\
2x_1-x_2+3x_3&=10.
}

\item \sistema{
9x_2-7x_3&=2,\\
-x_3&=-2,\\
-3x_1+6x_2+8x_3&=1.
}

\item \sistema{
-2x_1+x_2+6x_3&=18,\\
5x_1+8x_3&=-16,\\
3x_1+2x_2-10x_3&=-3.
}

\item \sistema{
3x_1+6x_2-6x_3&=9,\\
2x_1-5x_2+4x_3&=6,\\
5x_1+28x_2-26x_3&=-8.
}

\item \sistema{
3x_1+6x_2-6x_3&=9,\\
2x_1-5x_2+4x_3&=6,\\
-x_1+16x_2-14x_3&=-3.
}

\item \sistema{
-2x_1-6x_2-3x_3&=9,\\
-x_1+x_2-x_3&=1,\\
x_1-x_2+2x_3&=2.
}

\item \sistema{
x_1+x_2-x_3&=7,\\
4x_1-x_2+5x_3&=4,\\
2x_1+2x_2-3x_3&=0.
}

\item \sistema{
-x_1+x_3&=0,\\
x_2+3x_3&=1,\\
x_1-x_2&=-3.
}

\item \sistema{
x_1+x_2-x_3&=7,\\
4x_1-x_2+5x_3&=4,\\
6x_1+x_2+3x_3&=18.
}

\item \sistema{
x_1+2x_2-2x_3-x_4&=1,\\
-3x_1+4x_2+x_3-2x_4&=4,\\
-3x_1+14x_2+4x_3-7x_4&=3,\\
6x_1+12x_2-12x_3-6x_4&=5.
}

\item \sistema{
x_1-2x_2+3x_3&=0,\\
4x_1+x_2-x_3&=0,\\
2x_1-x_2+3x_3&=0.
}

\item \sistema{
x_1+x_2-x_3&=0,\\
4x_1-x_2+5x_3&=0,\\
6x_1+x_2+3x_3&=0.
}

\item \sistema{
x_1+2x_2-x_3&=4,\\
3x_1+4x_2-2x_3&=7.
}

\item \sistema{
x_1+2x_2-2x_3-x_4&=1,\\
-3x_1+4x_2+x_3-2x_4&=4,\\
-3x_1+14x_2-4x_3-7x_4&=3,\\
6x_1+12x_2-12x_3-6x_4&=5.
}

\item \sistema{
x_1+2x_2-4x_3&=4,\\
-2x_1-4x_2+8x_3&=-9.
}

\item \sistema{
x_1+2x_2-4x_3&=4,\\
-2x_1-4x_2+8x_3&=-8.
}

\item \sistema{
2x_1+6x_2-4x_3+2x_4&=4,\\
x_1-x_3+x_4&=5,\\
-3x_1+2x_2-2x_3&=-2.
}

\item \sistema{
x_1+2x_2-x_3+x_4&=7,\\
3x_1+6x_2-3x_3+3x_4&=21.
}

\item \sistema{
-x_1+2x_2-x_3+3x_4&=4,\\
-3x_1+6x_2-3x_3+9x_4&=12.
}

\item \sistema{
2x_1+x_2-x_3+x_4&=-2,\\
-3x_1+x_4&=1,\\
5x_2+8x_3&=3.
}

\item \sistema{
-2x_1+x_4&=1,\\
4x_2-x_3&=-1,\\
x_1+x_2&=-3.
}

\item \sistema{
x_1-2x_2+x_3+x_4&=2,\\
3x_1+2x_3-2x_4&=-8,\\
4x_2-x_3-x_4&=1,\\
5x_1+3x_3-x_4&=0.
}

\item \sistema{
x_1-2x_2+x_3+x_4&=2,\\
3x_1+2x_3-2x_4&=-8,\\
4x_2-x_3-x_4&=1,\\
5x_1+3x_3-x_4&=-3.
}

\item \sistema{
x_1+x_2&=4,\\
2x_1-3x_2&=7,\\
3x_1+2x_2&=8.
}

\item \sistema{
-2x_1+x_2&=0,\\
x_1+3x_2&=1,\\
3x_1-x_2&=-3.
}

\item \sistema{
x_1+x_2&=4,\\
2x_1-3x_2&=7,\\
3x_1-2x_2&=11.
}
\end{multicols}
\end{enumerate}
\end{Ejer}
%<<>>==<<>><<>>==<<>><<>>==<<>><<>>==<<>><<>>==<<>>
\begin{Ejer}
Forma escalonada y forma escalonada reducida. Determine si la matriz se encuentra en forma escalonada por renglones, en forma escalonada reducida por renglones o en ninguna de las dos.

\begin{multicols}{2}
\begin{enumerate}[resume]
\item $\begin{pmatrix}1&1&0\\0&1&0\\0&0&1\end{pmatrix}$
\item $\begin{pmatrix}3&5&2\\0&-2&5\\0&0&3\end{pmatrix}$
\item $\begin{pmatrix}2&0&0\\1&1&0\\0&0&1\end{pmatrix}$
\item $\begin{pmatrix}2&0&0\\0&1&0\\0&0&-1\end{pmatrix}$
\item $\begin{pmatrix}1&0&0&0\\0&1&0&0\\0&0&0&1\\0&0&0&0\end{pmatrix}$
\item $\begin{pmatrix}1&1&4&0\\0&0&1&3\\0&0&0&1\end{pmatrix}$
\item $\begin{pmatrix}0&1&0&0\\1&0&0&0\\0&0&0&0\end{pmatrix}$
\item $\begin{pmatrix}1&0&1&2\\0&1&3&4\end{pmatrix}$
\item $\begin{pmatrix}1&5&2\\0&1&5\end{pmatrix}$
\item $\begin{pmatrix}1&0\\0&1\\0&0\end{pmatrix}$
\item $\begin{pmatrix}1&0&0\\0&0&0\\0&0&1\end{pmatrix}$
\item $\begin{pmatrix}1&0&0&4\\0&1&0&5\\0&1&1&6\end{pmatrix}$
\end{enumerate}
\end{multicols}
\end{Ejer}
%<<>>==<<>><<>>==<<>><<>>==<<>><<>>==<<>><<>>==<<>>
\begin{Ejer} Reducción por renglones de matrices. Utilice operaciones elementales con renglones para obtener la forma escalonada y la forma escalonada reducida.

\begin{multicols}{2}
\begin{enumerate}[resume]
\item $\begin{pmatrix}1&1\\2&3\end{pmatrix}$
\item $\begin{pmatrix}-1&6\\4&2\end{pmatrix}$
\item $\begin{pmatrix}1&-1&1\\2&4&3\\5&6&-2\end{pmatrix}$
\item $\begin{pmatrix}1&-2&3\\-4&5&-6\\-1&1&1\end{pmatrix}$
\item $\begin{pmatrix}2&-4&8\\3&5&8\\-6&0&4\end{pmatrix}$
\item $\begin{pmatrix}2&-4&-2\\3&1&6\end{pmatrix}$
\item $\begin{pmatrix}3&-6&-3\\5&10&5\end{pmatrix}$
\item $\begin{pmatrix}2&-7\\3&5\\4&-3\end{pmatrix}$
\item $\begin{pmatrix}1&5&2\\-3&-14&-1\end{pmatrix}$
\end{enumerate}
\end{multicols}
\end{Ejer}
%<<>>==<<>><<>>==<<>><<>>==<<>><<>>==<<>><<>>==<<>>
\begin{Ejer}  Sistemas con parámetros y sistemas numéricos

\begin{enumerate}[resume]

\item Considere
\sistema{
5x_1+10x_2-20x_3&=a,\\
-6x_1-11x_2-21x_3&=b,\\
2x_1+4x_2+8x_3&=c.
}
Encuentre las condiciones sobre $a,b,c$ para que el sistema sea inconsistente.

\item Considere
\sistema{
2x_1-x_2+3x_3&=a,\\
3x_1+x_2-5x_3&=b,\\
-5x_1-5x_2+21x_3&=c.
}
Muestre que es inconsistente si $c\neq 2a-3b$.

\item Considere el sistema general
\sistema{
a_{11}x_1+a_{12}x_2+a_{13}x_3&=b_1,\\
a_{21}x_1+a_{22}x_2+a_{23}x_3&=b_2,\\
a_{31}x_1+a_{32}x_2+a_{33}x_3&=b_3.
}
Encuentre condiciones sobre los coeficientes $a_{ij}$ para que el sistema tenga solución única.

\item \sistema{
2x_2-x_3-4x_4&=2,\\
x_1-x_2+5x_3+2x_4&=-4,\\
3x_1+3x_2-7x_3-x_4&=4,\\
-x_1-2x_2+3x_3&=-7.
}

\item \sistema{
5.31x_1+1.14x_2+2.34x_3&=-3.2,\\
-6.44x_1-3.12x_2-1.97x_3&=1.1,\\
2.67x_1+4.32x_2+8.65x_3&=-2.4.
}

\item \sistema{
23.42x_1-16.89x_2+57.31x_3+82.6x_4&=2158.36,\\
-14.77x_1-38.29x_2+92.36x_3-4.36x_4&=-1123.02,\\
-77.21x_1+71.26x_2-16.55x_3+43.09x_4&=3248.71,\\
91.82x_1+81.43x_2+33.94x_3+57.22x_4&=235.25.
}

\item \sistema{
2.6x_1-4.3x_2+9.6x_3&=21.62,\\
-8.5x_1+3.6x_2+9.1x_3&=14.23,\\
12.3x_1-8.4x_2-0.6x_3&=12.61.
}

\item \sistema{
6.1x_1-2.4x_2+23.3x_3-16.4x_4-8.9x_5&=121.7,\\
-14.2x_1-31.6x_2-5.8x_3+9.6x_4+23.1x_5&=-87.7,\\
10.5x_1+46.1x_2-19.6x_3-8.8x_4-41.2x_5&=10.8,\\
37.3x_1-14.2x_2+62.0x_3+14.7x_4-9.6x_5&=61.3,\\
0.8x_1+17.7x_2-47.5x_3-50.2x_4+29.8x_5&=-27.8.
}

\item \sistema{
2.1x_1+4.2x_2-3.5x_3&=12.9,\\
-5.9x_1+2.7x_2+9.8x_3&=-1.6.
}

\item \sistema{
-13.6x_1+71.8x_2+46.3x_3&=-19.5,\\
41.3x_1-75.0x_2-82.9x_3&=46.4,\\
41.8x_1+65.4x_2-26.9x_3&=34.3.
}

\item \sistema{
-13.6x_1+71.8x_2+46.3x_3&=19.5,\\
41.3x_1-75.0x_2-82.9x_3&=46.4,\\
41.8x_1+65.4x_2-26.9x_3&=35.3.
}

\item \sistema{
5x_1-2x_2+11x_3-16x_4+12x_5&=105,\\
-6x_1+8x_2-14x_3-9x_4+26x_5&=-62,\\
7x_1-18x_2-12x_3+21x_4-2x_5&=53.
}

\item \sistema{
5x_1-2x_2+11x_3-16x_4+12x_5&=105,\\
-6x_1+8x_2-14x_3-9x_4+26x_5&=-62,\\
7x_1-18x_2-12x_3+21x_4-2x_5&=53,\\
-15x_1+42x_2+21x_3-17x_4+42x_5&=-63.
}

\item \sistema{
5x_1-2x_2+11x_3-16x_4+12x_5&=105,\\
-6x_1+8x_2-14x_3-9x_4+26x_5&=-62,\\
7x_1-18x_2-12x_3+21x_4-2x_5&=53,\\
-15x_1+42x_2+21x_3-17x_4+42x_5&=63.
}

\end{enumerate}
\end{Ejer}
%<<>>==<<>><<>>==<<>><<>>==<<>><<>>==<<>><<>>==<<>>
\begin{Ejer}  Encuentre todas las soluciones de los sistemas homogéneos.

\begin{enumerate}[resume]
\begin{multicols}{2}
\item \sistema{x_1-5x_2&=0,\\-x_1+5x_2&=0.}
\item \sistema{2x_1-x_2&=0,\\3x_1+4x_2&=0.}
\item \sistema{3x_1-5x_2&=0,\\5x_1+4x_2&=0,\\2x_1+5x_2&=0.}
\item \sistema{x_1-3x_2&=0,\\-2x_1+6x_2&=0.}
\item \sistema{x_1+x_2-x_3&=0,\\2x_1-4x_2+3x_3&=0,\\3x_1+7x_2-x_3&=0.}
\item \sistema{x_1+x_2-x_3&=0,\\2x_1-4x_2+3x_3&=0,\\-x_1-7x_2-6x_3&=0.}
\item \sistema{3x_1-5x_2+4x_3&=0,\\5x_1+4x_3&=0,\\2x_1+5x_2-2x_3&=0.}
\item \sistema{2x_1+3x_2-x_3&=0,\\6x_1-5x_2+7x_3&=0.}
\item \sistema{x_1-3x_2+2x_3&=0,\\3x_1+6x_2-3x_3&=0.}
\item \sistema{4x_1-x_2&=0,\\7x_1+3x_2&=0,\\-8x_1+6x_2&=0.}
\item \sistema{2x_1-3x_2-4x_3+5x_4&=0,\\7x_2+3x_4&=0.}
\item \sistema{2x_1-5x_2-6x_3-3x_4&=0,\\x_1+3x_2-5x_3+4x_4&=0.}
\item \sistema{x_1-2x_2+x_3+x_4&=0,\\3x_1+2x_3-2x_4&=0,\\4x_2-x_3-x_4&=0,\\5x_1+3x_3-x_4&=0.}
\item \sistema{-2x_1+7x_4&=0,\\x_1+2x_2-x_3+4x_4&=0,\\3x_1-x_3+5x_4&=0,\\4x_1+2x_2+3x_3&=0.}
\item \sistema{2x_1-x_2&=0,\\3x_1+5x_2&=0,\\7x_1-3x_2&=0,\\-2x_1+3x_2&=0.}
\item \sistema{x_1-3x_2&=0,\\-2x_1+6x_2&=0,\\4x_1-12x_2&=0.}
\item \sistema{-2x_1+6x_2&=0,\\x_1-3x_2&=0,\\-7x_1+21x_2&=0.}
\item \sistema{x_1+x_2-x_3&=0,\\4x_1-x_2+5x_3&=0,\\-2x_1+x_2-2x_3&=0,\\3x_1+2x_2-6x_3&=0.}
\item \sistema{3x_1-5x_2+12x_3+10x_4&=0,\\5x_1+4x_2+20x_3-8x_4&=0,\\2x_1+5x_2+8x_3-10x_4&=0.}
\item \sistema{4x_1+10x_2-6x_3&=0,\\-6x_1-9x_2-9x_3&=0,\\-x_1+2x_2-12x_3&=0,\\-x_1-6x_2-12x_3&=0.}
\end{multicols}
\end{enumerate}
\end{Ejer}
%<<>>==<<>><<>>==<<>><<>>==<<>><<>>==<<>><<>>==<<>>
\begin{Ejer} Ejercicios conceptuales y con parámetros

\begin{enumerate}[resume]
\item Muestre que
\sistema{a_{11}x_1+a_{12}x_2&=0,\\a_{21}x_1+a_{22}x_2&=0}
tiene infinitas soluciones si y sólo si
\[
a_{11}a_{22}-a_{12}a_{21}=0.
\]

\item Considere
\sistema{
2x_1-3x_2+5x_3&=0,\\
-x_1+7x_2-x_3&=0,\\
4x_1-11x_2+kx_3&=0.
}
¿Para qué valor de $k$ tendrá soluciones no triviales?

\item Considere el sistema homogéneo general $3\times3$
\sistema{
a_{11}x_1+a_{12}x_2+a_{13}x_3&=0,\\
a_{21}x_1+a_{22}x_2+a_{23}x_3&=0,\\
a_{31}x_1+a_{32}x_2+a_{33}x_3&=0.
}
Encuentre condiciones sobre los coeficientes $a_{ij}$ tales que la solución trivial sea la única.

\item Determine para qué valores de $K$ el sistema tiene solución única:
\sistema{
Kx+y+z&=1,\\
x+Ky+z&=1,\\
x+y+Kz&=1.
}

\item Para el sistema
\sistema{
2x-y-Kz&=0,\\
x-y-2z&=1,\\
-x+2z&=K,
}
determine para qué valores de $K$:
\begin{enumerate}[label=\alph*)]
\item no tiene solución;
\item tiene infinitas soluciones;
\item tiene solución única.
\end{enumerate}

\item \sistema{
4.23x_1+10.28x_2-6.36x_3&=0,\\
3.28x_1-5.39x_2+4.25x_3&=0.
}

\item \sistema{
-13.6x_1+71.8x_2+46.3x_3&=0,\\
41.3x_1-75.0x_2-82.9x_3&=0,\\
41.8x_1+65.4x_2-26.9x_3&=0.
}

\item \sistema{
2.1x_1+4.2x_2-3.5x_3&=0,\\
-5.9x_1+2.7x_2+8.9x_3&=0.
}

\item \sistema{
5x_1-2x_2+11x_3-16x_4+12x_5&=0,\\
-6x_1+8x_2-14x_3-9x_4+26x_5&=0,\\
7x_1-18x_2-12x_3+21x_4-2x_5&=0,\\
-x_1+11x_2-9x_3+13x_4-20x_5&=0.
}
\end{enumerate}
\end{Ejer}
%<<>>==<<>><<>>==<<>><<>>==<<>><<>>==<<>><<>>==<<>>
\begin{Ejer} Resuelva mediante eliminación de Gauss--Jordan. No utilice regla de Cramer.

\begin{enumerate}[resume]
\begin{multicols}{3}
\item 
\sistema{
x+2y&=7,\\
3x-y&=7.
}

\item 
\sistema{
2x-3y&=-4,\\
5x+2y&=13.
}

\item 
\sistema{
4x+6y&=10,\\
2x+3y&=5.
}

\item 
\sistema{
3x-2y&=8,\\
-6x+4y&=-13.
}

\item 
\sistema{
x+y+z&=6,\\
2x-y+3z&=9,\\
3x+2y-z&=4.
}

\item
\sistema{
2x-y+z&=5,\\
x+3y-2z&=-1,\\
3x+2y+4z&=14.
}

\item
\sistema{
y+2z&=5,\\
2x-y+z&=1,\\
x+3y-z&=6.
}

\item 
\sistema{
x-2y+3z&=4,\\
2x-4y+6z&=8,\\
3x+y-z&=5.
}

\item 
\sistema{
x+2y-z&=3,\\
2x+4y-2z&=7,\\
-x+y+3z&=2.
}

\item 
\sistema{
2x+y-z&=3,\\
-4x-2y+2z&=-6,\\
x-3y+4z&=5.
}

\item 
\sistema{
3x-2y+4z&=7,\\
2x+5y-z&=-3,\\
-x+3y+2z&=8.
}

\item
\sistema{
4x-y+2z&=9,\\
-2x+3y+5z&=-1,\\
6x+2y-z&=12.
}

\item
\sistema{
2x-5y+3z&=1,\\
4x-10y+6z&=2,\\
-x+4y-2z&=3.
}

\item
\sistema{
x+3y-2z&=4,\\
2x+6y-4z&=8,\\
-3x-9y+6z&=-12.
}

\item 
\sistema{
x+y+2z&=4,\\
2x+3y+5z&=9,\\
3x+4y+7z&=14.
}
\end{multicols}
\end{enumerate}
\end{Ejer}
%<<>>==<<>><<>>==<<>><<>>==<<>><<>>==<<>><<>>==<<>>

\begin{Ejer} 

\begin{enumerate}[resume]
\begin{multicols}{2}
\item
\sistema{
x_1+x_2-x_3+x_4&=2,\\
2x_1-x_2+x_3+2x_4&=7,\\
-x_1+2x_2+2x_3-x_4&=-1,\\
3x_1+x_2+x_3+x_4&=8.
}

\item
\sistema{
x_1-2x_2+x_3+3x_4&=7,\\
2x_1+x_2-3x_3+x_4&=-2,\\
3x_1-x_2+2x_3-2x_4&=5,\\
-x_1+4x_2+x_3+x_4&=1.
}

\item
\sistema{
2x_1+x_2-x_3+x_4&=3,\\
4x_1+2x_2-2x_3+2x_4&=6,\\
x_1-3x_2+2x_3-x_4&=-4,\\
3x_1-2x_2+x_3&=-1.
}

\item 
\sistema{
x_1+2x_2-x_3+x_4&=5,\\
2x_1+4x_2-2x_3+2x_4&=10,\\
-x_1+x_2+3x_3-2x_4&=1,\\
3x_1+7x_2+x_3&=12.
}

\item
\sistema{
x_2+2x_3-x_4&=4,\\
x_1-x_2+x_3+2x_4&=1,\\
2x_1+x_2-3x_3+x_4&=6,\\
-x_1+3x_2+x_3-x_4&=-2.
}

\item 
\sistema{
3x_1-2x_2+x_3+4x_4&=11,\\
-2x_1+5x_2+3x_3-x_4&=-4,\\
x_1+4x_2-2x_3+2x_4&=7,\\
5x_1+x_2+4x_3-3x_4&=6.
}

\item
\sistema{
2x_1-3x_2+4x_3-x_4&=5,\\
4x_1-6x_2+8x_3-2x_4&=10,\\
-x_1+2x_2-x_3+3x_4&=-2,\\
3x_1-x_2+5x_3+2x_4&=8.
}

\item
\sistema{
x_1+2x_2+x_3-x_4&=3,\\
2x_1+4x_2+2x_3-2x_4&=6,\\
-x_1+x_2+3x_3+2x_4&=5,\\
3x_1+7x_2+5x_3&=12.
}
\end{multicols}
\end{enumerate}
\end{Ejer}

% =========================================================
\begin{Ejer} Gauss--Jordan: sistemas $5\times5$. Estos problemas están pensados para ilustrar exhaustivamente la técnica. Se recomienda trabajarlos por etapas, identificando el pivote actual y la columna que se desea limpiar.

\begin{enumerate}[resume]
\begin{multicols}{2}
\item 
\sistema{
x_1+x_2-x_3+2x_4+x_5&=6,\\
2x_1-x_2+x_3+x_4-2x_5&=1,\\
-x_1+2x_2+2x_3-x_4+x_5&=3,\\
3x_1+x_2+x_3+2x_4+x_5&=10,\\
2x_1+3x_2-x_3+x_4+2x_5&=9.
}

\item
\sistema{
x_1-2x_2+x_3+x_4-3x_5&=4,\\
2x_1+x_2-2x_3+3x_4+x_5&=-1,\\
-x_1+3x_2+x_3-2x_4+2x_5&=5,\\
3x_1-x_2+4x_3+x_4-x_5&=7,\\
2x_1+2x_2-x_3-3x_4+4x_5&=2.
}

\item 
\sistema{
2x_1-x_2+3x_3+x_4-2x_5&=7,\\
-x_1+4x_2+x_3-3x_4+x_5&=-2,\\
3x_1+2x_2-2x_3+x_4+4x_5&=9,\\
x_1-3x_2+5x_3+2x_4-x_5&=4,\\
4x_1+x_2+x_3-2x_4+3x_5&=11.
}

\item
\sistema{
x_1+x_2-x_3+x_4+x_5&=2,\\
2x_1+2x_2-2x_3+2x_4+2x_5&=4,\\
x_1-x_2+2x_3-3x_4+x_5&=1,\\
3x_1+x_2+4x_3-5x_4+3x_5&=5,\\
2x_1+4x_2-5x_3+7x_4+x_5&=6.
}

\item
\sistema{
x_1-2x_2+3x_3-x_4+2x_5&=5,\\
2x_1-4x_2+6x_3-2x_4+4x_5&=10,\\
-x_1+x_2+x_3+3x_4-x_5&=2,\\
3x_1-5x_2+7x_3+x_4+3x_5&=12,\\
2x_1-3x_2+4x_3+4x_4+2x_5&=9.
}

\item 
\sistema{
3x_1-2x_2+x_3+4x_4-x_5&=8,\\
-2x_1+5x_2+3x_3-x_4+2x_5&=-3,\\
x_1+4x_2-2x_3+2x_4+3x_5&=6,\\
5x_1+x_2+4x_3-3x_4+x_5&=10,\\
-3x_1+2x_2+5x_3+x_4-4x_5&=-1.
}
\end{multicols}
\end{enumerate}
\end{Ejer}
% =========================================================

\begin{Ejer} Forma escalonada y forma escalonada reducida:

Para cada matriz indique si está en:
\begin{enumerate}[label=\alph*)]
\item forma escalonada por renglones;
\item forma escalonada reducida por renglones;
\item ninguna de las anteriores.
\end{enumerate}
Justifique su respuesta indicando pivotes y condiciones que se cumplen o se violan.

\begin{multicols}{2}
\begin{enumerate}[resume]
\item $\begin{pmatrix}1&2&0&3\\0&1&4&-1\\0&0&1&2\end{pmatrix}$
\item $\begin{pmatrix}1&0&3&0\\0&1&-2&0\\0&0&0&1\end{pmatrix}$
\item $\begin{pmatrix}1&2&0\\0&0&1\\0&0&0\end{pmatrix}$
\item $\begin{pmatrix}0&1&2\\1&0&3\\0&0&0\end{pmatrix}$
\item $\begin{pmatrix}1&0&2&4\\0&1&-3&1\\0&0&0&0\end{pmatrix}$
\item $\begin{pmatrix}2&1&0\\0&3&4\\0&0&5\end{pmatrix}$
\item $\begin{pmatrix}1&2&0\\0&0&0\\0&1&3\end{pmatrix}$
\item $\begin{pmatrix}1&0&0&2\\0&1&0&-1\\0&0&1&4\end{pmatrix}$
\item $\begin{pmatrix}1&3&0&2\\0&1&0&4\\0&0&0&0\end{pmatrix}$
\item $\begin{pmatrix}0&0&1&2\\0&0&0&1\\0&0&0&0\end{pmatrix}$
\end{enumerate}
\end{multicols}
\end{Ejer}

% =========================================================
\begin{Ejer} Reducción de matrices por renglones. Obtenga primero una forma escalonada y después la forma escalonada reducida.

\begin{multicols}{2}
\begin{enumerate}[resume]
\item $\begin{pmatrix}2&1\\5&3\end{pmatrix}$
\item $\begin{pmatrix}3&-2\\-6&5\end{pmatrix}$
\item $\begin{pmatrix}1&2&-1\\2&5&1\\-1&1&4\end{pmatrix}$
\item $\begin{pmatrix}2&-1&3\\4&1&7\\-2&3&-1\end{pmatrix}$
\item $\begin{pmatrix}1&3&-2\\2&6&-4\\-1&2&5\end{pmatrix}$
\item $\begin{pmatrix}2&-4&6&1\\1&-2&3&4\\3&1&-2&0\end{pmatrix}$
\item $\begin{pmatrix}1&2&0&-1\\2&4&1&3\\-1&-2&2&5\end{pmatrix}$
\item $\begin{pmatrix}0&2&-1\\3&1&4\\6&2&8\end{pmatrix}$
\end{enumerate}
\end{multicols}
\end{Ejer}
% =========================================================
\begin{Ejer} Sistemas con parámetros

Analice cada sistema mediante reducción por renglones. Determine los valores del parámetro
para los cuales el sistema tiene solución única, infinitas soluciones o ninguna solución.

\begin{enumerate}[resume]

\item
\sistema{
x+y+z&=2,\\
2x+3y+4z&=5,\\
x+2y+kz&=3.
}

\item
\sistema{
kx+y&=1,\\
x+ky&=1.
}

\item
\sistema{
x+2y-z&=1,\\
2x+4y-2z&=k,\\
x-y+3z&=4.
}

\item 
\sistema{
x+y+z&=a,\\
2x-y+z&=b,\\
3x+2z&=c.
}
Encuentre una condición entre $a,b,c$ que garantice la consistencia si los renglones
de la matriz de coeficientes resultan dependientes.

\item
\sistema{
(k-1)x+y+z&=1,\\
x+(k-1)y+z&=1,\\
x+y+(k-1)z&=1.
}

\item 
\sistema{
x+2y-z+w&=1,\\
2x+4y-2z+2w&=k,\\
-x+y+3z-w&=2,\\
3x+5y+z+w&=4.
}

\end{enumerate}
\end{Ejer}
% =========================================================
\begin{Ejer} Sistemas homogéneos. Encuentre todas las soluciones. Determine si existen soluciones no triviales y describa el conjunto solución mediante parámetros.

\begin{enumerate}[resume]
\begin{multicols}{2}
\item 
\sistema{x-4y&=0,\\2x-8y&=0.}

\item 
\sistema{3x-y&=0,\\2x+5y&=0.}

\item 
\sistema{x+y-z&=0,\\2x-y+3z&=0.}

\item 
\sistema{x-2y+z&=0,\\2x-4y+2z&=0,\\-3x+6y-3z&=0.}

\item 
\sistema{
2x_1-x_2+x_3&=0,\\
x_1+3x_2-2x_3&=0,\\
3x_1+2x_2-x_3&=0.
}

\item 
\sistema{
x_1+2x_2-x_3+x_4&=0,\\
2x_1+4x_2-2x_3+2x_4&=0,\\
-x_1+x_2+3x_3-x_4&=0.
}

\item 
\sistema{
2x_1-x_2+3x_3-x_4&=0,\\
-x_1+4x_2+x_3+2x_4&=0,\\
3x_1+2x_2-2x_3+x_4&=0.
}

\item 
\sistema{
x_1-2x_2+x_3+x_4-x_5&=0,\\
2x_1-4x_2+2x_3+2x_4-2x_5&=0,\\
-x_1+x_2+3x_3-2x_4+x_5&=0,\\
3x_1-5x_2+5x_3+x_4-x_5&=0.
}
\end{multicols}
\end{enumerate}
\end{Ejer}
% =========================================================
\begin{Ejer} Problemas conceptuales

\begin{enumerate}[resume]

\item Explique por qué un sistema homogéneo siempre es consistente.

\item ¿Puede un sistema homogéneo tener exactamente dos soluciones? Justifique.

\item Construya un sistema $3\times3$ cuya forma escalonada reducida tenga exactamente
dos pivotes. Describa su conjunto solución.

\item Construya un sistema de tres ecuaciones con tres incógnitas que sea incompatible,
pero que se vuelva compatible al cambiar solamente un término independiente.

\item Construya dos matrices aumentadas distintas que sean equivalentes por renglones
y representen sistemas con el mismo conjunto solución.

\item Si una matriz aumentada contiene el renglón
\[
\begin{pmatrix}0&0&0&0&|&5\end{pmatrix},
\]
¿qué puede concluirse inmediatamente? Explique la interpretación algebraica.

\item Si un sistema de cinco incógnitas tiene tres pivotes y es consistente, ¿cuántas
variables libres tiene? ¿Qué puede decirse sobre el número de soluciones?

\item Compare eliminación de Gauss y Gauss--Jordan. ¿En qué punto puede detenerse
Gauss y qué trabajo adicional realiza Gauss--Jordan?

\end{enumerate}
\end{Ejer}


\begin{Ejer}Resolver y clasificar

\begin{enumerate}
\begin{multicols}{2}
\item 
\[
\left\{\begin{aligned}
2x+y-3z&=-9\\
3x+y&=-4\\
\end{aligned}\right.
\]
\item 
\[
\left\{\begin{aligned}
3x-y&=11\\
2x+4y-3z&=21\\
\end{aligned}\right.
\]
\item 
\[
\left\{\begin{aligned}
-3x-2y+z&=12\\
-3x-4y-2z&=18\\
\end{aligned}\right.
\]
\item 
\[
\left\{\begin{aligned}
x+3y-z&=12\\
3x+2y+2z&=20\\
\end{aligned}\right.
\]
\item 
\[
\left\{\begin{aligned}
-4x-y+3z&=1\\
x-4y&=-13\\
\end{aligned}\right.
\]
\item 
\[
\left\{\begin{aligned}
x+4y-z&=14\\
-4x-y&=1\\
\end{aligned}\right.
\]
\item 
\[
\left\{\begin{aligned}
3x-2y+2z&=10\\
3x-3y-2z&=-3\\
\end{aligned}\right.
\]
\item 
\[
\left\{\begin{aligned}
-4x+4y-4z&=4\\
-2x-y-4z&=10\\
\end{aligned}\right.
\]
\item 
\[
\left\{\begin{aligned}
2x-3y+z&=-17\\
2x+3y&=8\\
\end{aligned}\right.
\]
\item 
\[
\left\{\begin{aligned}
-2x+4y-2z&=10\\
2x+2y&=0\\
\end{aligned}\right.
\]
\item 
\[
\left\{\begin{aligned}
3x+2y+4z&=23\\
x+2y-4z&=-7\\
\end{aligned}\right.
\]
\item 
\[
\left\{\begin{aligned}
-4x-z&=0\\
2x-y-3z&=-18\\
\end{aligned}\right.
\]
\end{multicols}
\end{enumerate}
\end{Ejer}

\begin{Ejer}Resolver y clasificar

\begin{enumerate}
\begin{multicols}{2}
\item 
\[
\left\{\begin{aligned}
-3x+2y&=9\\
-4x+2y&=10\\
-2x+4y&=14\\
\end{aligned}\right.
\]
\item
\[
\left\{\begin{aligned}
3x-3y&=-21\\
y&=4\\
2x&=-6\\
\end{aligned}\right.
\]
\item 
\[
\left\{\begin{aligned}
x-y&=-4\\
2y&=6\\
x+2y&=5\\
\end{aligned}\right.
\]
\item
\[
\left\{\begin{aligned}
-4x-3y&=-24\\
-2x+y&=-2\\
4x+3y&=24\\
\end{aligned}\right.
\]
\item
\[
\left\{\begin{aligned}
2y&=4\\
3x-2y&=2\\
x-2y&=-2\\
\end{aligned}\right.
\]
\item 
\[
\left\{\begin{aligned}
2x+2y&=6\\
-2x+4y&=-6\\
4x+4y&=13\\
\end{aligned}\right.
\]
\item 
\[
\left\{\begin{aligned}
-4x-3y&=5\\
2x-2y&=8\\
-x+y&=-4\\
\end{aligned}\right.
\]
\item
\[
\left\{\begin{aligned}
-x+4y&=4\\
4x+2y&=20\\
x-4y&=-4\\
\end{aligned}\right.
\]
\item 
\[
\left\{\begin{aligned}
4y&=16\\
-4x+3y&=8\\
3y&=12\\
\end{aligned}\right.
\]
\item
\[
\left\{\begin{aligned}
-4x+y&=-3\\
-4x+y&=-3\\
2x+2y&=4\\
\end{aligned}\right.
\]
\item 
\[
\left\{\begin{aligned}
3x+y&=15\\
x-2y&=-2\\
-4x-4y&=-28\\
\end{aligned}\right.
\]
\item 
\[
\left\{\begin{aligned}
-3x+y&=-2\\
-x+2y&=6\\
-6x+2y&=-3\\
\end{aligned}\right.
\]

\end{multicols}
\end{enumerate}
\end{Ejer}

\begin{Ejer}Resolver y clasificar

\begin{enumerate}
\begin{multicols}{2}
\item 
\[
\left\{\begin{aligned}
4x-3y+4z+2w&=-12\\
3x-4y&=-2\\
3x+y-3z+w&=11\\
\end{aligned}\right.
\]
\item 
\[
\left\{\begin{aligned}
x+y+3z-2w&=19\\
-2x-2y+3z-4w&=22\\
4x-4z-w&=-17\\
\end{aligned}\right.
\]
\item 
\[
\left\{\begin{aligned}
-4x+2y+z&=16\\
3x+y+z-4w&=-24\\
-x+3y+4z+w&=8\\
\end{aligned}\right.
\]
\item 
\[
\left\{\begin{aligned}
3x-y-3z-4w&=9\\
4x+3y+3z+2w&=6\\
-3x+y+3z+4w&=-9\\
\end{aligned}\right.
\]
\item
\[
\left\{\begin{aligned}
-4x+4y+2z+3w&=13\\
x-y+3z+3w&=11\\
-4x+4y+4z-3w&=13\\
\end{aligned}\right.
\]
\item 
\[
\left\{\begin{aligned}
-3x+2y-4z+2w&=-6\\
-x-4y+3z-w&=7\\
-6x+4y-8z+4w&=-11\\
\end{aligned}\right.
\]
\item 
\[
\left\{\begin{aligned}
4x+y-3w&=-2\\
3x+4y-4z&=-9\\
-3x+2y+2w&=1\\
\end{aligned}\right.
\]
\item 
\[
\left\{\begin{aligned}
-x-3y+3z-2w&=-5\\
x-2y-w&=7\\
x+3y-3z+2w&=5\\
\end{aligned}\right.
\]
\item 
\[
\left\{\begin{aligned}
4x+3y+4z&=24\\
3x-2y-4w&=7\\
3x-2y-2w&=1\\
\end{aligned}\right.
\]
\item 
\[
\left\{\begin{aligned}
3x+2y-4z+2w&=18\\
4x+y+4z-w&=-2\\
-x-4y+4z-2w&=-14\\
\end{aligned}\right.
\]
\item 
\[
\left\{\begin{aligned}
-2x-3y+3z+3w&=5\\
-x+y+2w&=1\\
4x+4y-3z-2w&=0\\
\end{aligned}\right.
\]
\item 
\[
\left\{\begin{aligned}
x+3y-2w&=5\\
2x-3z+3w&=6\\
2x+6y-4w&=11\\
\end{aligned}\right.
\]

\end{multicols}
\end{enumerate}
\end{Ejer}

\begin{Ejer}Resolver y clasificar

\begin{enumerate}
\begin{multicols}{2}

\item 
\[
\left\{\begin{aligned}
-3x-3y+2z+2w+u&=-4\\
2y+4z-2w-u&=17\\
4x+4y+3z-3w+4u&=33\\
\end{aligned}\right.
\]
\item 
\[
\left\{\begin{aligned}
-4x-y+z+w-4u&=6\\
-3w+u&=1\\
4x-2y+2w+u&=-6\\
\end{aligned}\right.
\]
\item 
\[
\left\{\begin{aligned}
-4x-y+2z+3w+4u&=5\\
2x+y-4z-2w-4u&=-12\\
-4x-y+3z-2w+4u&=13\\
\end{aligned}\right.
\]
\item 
\[
\left\{\begin{aligned}
-3x+z+3w-3u&=6\\
-2x+y-3z-2w&=19\\
3x-z-3w+3u&=-6\\
\end{aligned}\right.
\]
\item 
\[
\left\{\begin{aligned}
-2x+2z+3w+3u&=7\\
-4x-4y+3z+4w+4u&=13\\
4x-y+2u&=-8\\
\end{aligned}\right.
\]
\item
\[
\left\{\begin{aligned}
3y-z+w+2u&=5\\
x-4y-4z+4w-u&=-2\\
6y-2z+2w+4u&=11\\
\end{aligned}\right.
\]
\item 
\[
\left\{\begin{aligned}
-4x-3y+4z+3w+4u&=20\\
-3x-y-3z+3w-u&=-14\\
-3x+y+z+w+4u&=0\\
\end{aligned}\right.
\]
\item 
\[
\left\{\begin{aligned}
-4x+2y+z-u&=14\\
4y+3u&=-6\\
4x-2y-z+u&=-14\\
\end{aligned}\right.
\]
\item 
\[
\left\{\begin{aligned}
x-4y-4z-w-2u&=23\\
3x-y-2z-u&=8\\
-4x-3y-u&=8\\
\end{aligned}\right.
\]
\item 
\[
\left\{\begin{aligned}
-4x+2y+2z-3w-2u&=1\\
-4x+y-2z-4w-2u&=-13\\
-3x-3z-w+2u&=-15\\
\end{aligned}\right.
\]
\item 
\[
\left\{\begin{aligned}
3x-2y+4z+w&=13\\
x-4y-3z+2w+u&=6\\
4x-3y-3z+3w+2u&=1\\
\end{aligned}\right.
\]
\item 
\[
\left\{\begin{aligned}
-4x-3z+w+2u&=6\\
3x-4y+4z-3w&=9\\
-8x-6z+2w+4u&=13\\
\end{aligned}\right.
\]

\end{multicols}
\end{enumerate}
\end{Ejer}

\begin{Ejer}Resolver y clasificar

\begin{enumerate}
\begin{multicols}{2}

\item 
\[
\left\{\begin{aligned}
3x+4y-2z+3w+4u&=34\\
2x+4y-4z+2w+4u&=28\\
4x+2y-4z+3w&=10\\
-2x-4y-3w+u&=-14\\
\end{aligned}\right.
\]
\item
\[
\left\{\begin{aligned}
-3x-y-z-4w-u&=-17\\
-2x+4y+4z-w&=24\\
2x-3y-3z-w+4u&=-25\\
x-2z-3w+3u&=-14\\
\end{aligned}\right.
\]
\item 
\[
\left\{\begin{aligned}
2x-4y+3z-w+3u&=-3\\
4x-2y-2z+4w-4u&=0\\
x+4y-3z-3w+3u&=24\\
-3x-3y+3z-3w+3u&=-15\\
\end{aligned}\right.
\]
\item 
\[
\left\{\begin{aligned}
-3y-2z+3w-3u&=20\\
4x+3y+2z+4w-u&=16\\
4x+2y+z-3w+2u&=-21\\
3y+2z-3w+3u&=-20\\
\end{aligned}\right.
\]
\item 
\[
\left\{\begin{aligned}
-4x+y-3z-2w&=12\\
x+3y+2z+4w-u&=-11\\
-x+2y+2z-2u&=7\\
3x-3w-2u&=0\\
\end{aligned}\right.
\]
\item 
\[
\left\{\begin{aligned}
-2x+2y+2z-2w&=16\\
-2x-2y-2z+3u&=-7\\
-y+3z-3u&=10\\
-4x+4y+4z-4w&=33\\
\end{aligned}\right.
\]
\item 
\[
\left\{\begin{aligned}
4x-3y-2w+u&=23\\
x+y+z+4w+4u&=10\\
2x+y-4w+4u&=37\\
-2x+4y+z-2w-4u&=-13\\
\end{aligned}\right.
\]
\item 
\[
\left\{\begin{aligned}
-4x-2y+4z-4w+3u&=-1\\
4x+3y+z+3w&=2\\
-3x-y-2z+2w+3u&=16\\
4x+2y-4z+4w-3u&=1\\
\end{aligned}\right.
\]
\item 
\[
\left\{\begin{aligned}
-3x-y-4z-4w-3u&=-7\\
x+2y+3w-u&=8\\
-2x-4y+2w-3u&=15\\
3x+4y-2w-2u&=-3\\
\end{aligned}\right.
\]
\item 
\[
\left\{\begin{aligned}
-3x+y+2z+w-u&=-4\\
-4x-2y-z-4w&=-18\\
-x+3y+z+2w+u&=-1\\
-2x+2y+z+3w-2u&=5\\
\end{aligned}\right.
\]
\item 
\[
\left\{\begin{aligned}
3x-2y+3z+2w+2u&=24\\
-2x+4y+4z-4w&=10\\
-2x-w&=-4\\
3x+2y-3z-2w-u&=-2\\
\end{aligned}\right.
\]
\item 
\[
\left\{\begin{aligned}
-3x-2y-4z-w-4u&=-14\\
-2x-y-4z+4w&=-13\\
-x-y-2z-4w+4u&=28\\
-6x-4y-8z-2w-8u&=-27\\
\end{aligned}\right.
\]
\end{multicols}
\end{enumerate}
\end{Ejer}

\begin{Ejer}Resolver y clasificar

\begin{enumerate}
\begin{multicols}{2}
\item
\[
\left\{\begin{aligned}
x-4y-4z-2u+3v&=-16\\
-x+y-4z+4w+3u-4v&=-10\\
2y-z-4w+4u-v&=12\\
-x+2y+3z+2w-2u-3v&=7\\
2x+4y+3w+4u+v&=-18\\
\end{aligned}\right.
\]
\item 
\[
\left\{\begin{aligned}
x-2y-z-2w-3u-3v&=1\\
-4x+y+w+3u-2v&=-18\\
-2x-y+z+3w-3u+2v&=7\\
-4x+y-3z+3w-3v&=-9\\
x-3y+2z+3w-3u&=18\\
\end{aligned}\right.
\]
\item 
\[
\left\{\begin{aligned}
2y-3z+w+3u-4v&=-7\\
x+2y+4w-4u+2v&=-15\\
x+3y-4z+w+3u-2v&=-9\\
4x-4y+4z+3w-2u+4v&=-17\\
2x-4y+3z+3w-v&=-16\\
\end{aligned}\right.
\]
\item 
\[
\left\{\begin{aligned}
3x+2y+4z-w-u+2v&=-13\\
-x+4y-3z-2w+2u-2v&=1\\
-3x+y-3z+3w-4u+3v&=-7\\
4x+y+3z+4w-3u-2v&=2\\
-3x-2y-4z+w+u-2v&=13\\
\end{aligned}\right.
\]
\item 
\[
\left\{\begin{aligned}
-4x-4y-4z-w-u+4v&=-29\\
3z+2w+2u+4v&=0\\
-4x+y+z+w+3v&=-23\\
3x-2y-2z+3w-u-v&=3\\
-2x-3y-z+3w-2u+v&=-21\\
\end{aligned}\right.
\]
\item 
\[
\left\{\begin{aligned}
-4x-3y-z+3w-u+3v&=23\\
-3x-2z+2w&=9\\
-2x+y+3z+4w-3v&=11\\
-2x+3y+3z-w+3u-2v&=-7\\
-8x-6y-2z+6w-2u+6v&=47\\
\end{aligned}\right.
\]
\item 
\[
\left\{\begin{aligned}
3y-3z-2w-2u-4v&=5\\
3x-3y-4z+3w-4u-v&=-18\\
3x-y+2z+2w+u+2v&=6\\
4x+z+3w-2u-3v&=2\\
x-3z+3u-v&=-4\\
\end{aligned}\right.
\]
\item 
\[
\left\{\begin{aligned}
2x-2y-2z+w+4v&=-12\\
-3x+3y-4w+v&=13\\
3x-3y-z+w+4u&=-9\\
4x+3y-2z+2u+3v&=8\\
-2x+2y+2z-w-4v&=12\\
\end{aligned}\right.
\]
\item 
\[
\left\{\begin{aligned}
x+3y-2w+3u+2v&=-28\\
-x-4y-3z+w+3u-2v&=9\\
-2x+y-3w+4u+2v&=-25\\
4x+2z-4w+4u-4v&=-16\\
-2x-3y-4z-3w-u-4v&=1\\
\end{aligned}\right.
\]
\item 
\[
\left\{\begin{aligned}
-4y-3z-w-v&=8\\
3x-4z-4w-3u-3v&=-20\\
-x+4y+w+2u+3v&=3\\
-x-2z-w-u+4v&=12\\
4x-3y+w-4u-3v&=-4\\
\end{aligned}\right.
\]
\item 
\[
\left\{\begin{aligned}
-3x-y-4z-3u-4v&=-32\\
-x+y-3z-2w-u-2v&=-25\\
3x-y-3z+3w-4u+2v&=19\\
-4x+4y+z+4w+3u-4v&=-5\\
2x+2y-4z-2w+4u&=-12\\
\end{aligned}\right.
\]
\item 
\[
\left\{\begin{aligned}
-2x-y-3z-2w-2u-2v&=-5\\
-3x-y-3z+2w-u+4v&=-10\\
-3x+4y+w-3u&=-23\\
2x-2y+z-4w-4u-3v&=4\\
-4x-2y-6z-4w-4u-4v&=-9\\
\end{aligned}\right.
\]
\end{multicols}
\end{enumerate}
\end{Ejer}

\begin{Ejer}Resolver y clasificar

\begin{enumerate}
\begin{multicols}{2}
\item

\[
\left\{\begin{aligned}
3x-3y+4z-2w+u+4v&=-13\\
-x+2y+3z-3w-3v&=16\\
2y+2z+3w+u-2v&=19\\
\end{aligned}\right.
\]
\item
\[
\left\{\begin{aligned}
-3x+2y-3z-w-u-2v&=25\\
-y-3z-2w+4u+3v&=3\\
-2x+4y+2w+2u+2v&=16\\
\end{aligned}\right.
\]
\item 
\[
\left\{\begin{aligned}
4x-y+2z+4w+4u-4v&=-4\\
x-y+2z+2w-4u-3v&=13\\
4x-y+4z-w-2u-v&=2\\
\end{aligned}\right.
\]
\item 
\[
\left\{\begin{aligned}
2y+z+w-2v&=5\\
4x+3y-3w-2u-3v&=23\\
-2y-z-w+2v&=-5\\
\end{aligned}\right.
\]
\item 
\[
\left\{\begin{aligned}
-2x+3y-3z+2w-u+2v&=-20\\
x+2z+3w+2u&=5\\
-x-y-3z-u-4v&=-4\\
\end{aligned}\right.
\]
\item
\[
\left\{\begin{aligned}
-y+z-3w-2u+3v&=19\\
4x+3y-3z-3w+4u+v&=-17\\
-2y+2z-6w-4u+6v&=39\\
\end{aligned}\right.
\]
\item 
\[
\left\{\begin{aligned}
-2x+y+3z+2w+2u-3v&=-26\\
-3z+w-4u&=9\\
-x-4y+z+w-2v&=-26\\
\end{aligned}\right.
\]
\item 
\[
\left\{\begin{aligned}
3x-4y+4z+2w-4u&=-19\\
-4x-3y-3z+w-2u+3v&=6\\
-3x+4y-4z-2w+4u&=19\\
\end{aligned}\right.
\]
\item 
\[
\left\{\begin{aligned}
-2x-2y+4z-4w+3u-2v&=4\\
-3x+3y+3z-2w-4u-2v&=3\\
x+4y+3z+w-4u+4v&=0\\
\end{aligned}\right.
\]
\item 
\[
\left\{\begin{aligned}
4x-3y-z+w-2u-2v&=29\\
x-2y+3z+w+3u+2v&=-3\\
x+3y+4z-2w+2u+4v&=-15\\
\end{aligned}\right.
\]
\item 
\[
\left\{\begin{aligned}
-x+2z-2w-4u&=-8\\
3x+y-3z-2w-u-v&=6\\
-4x-y+w-u+v&=-1\\
\end{aligned}\right.
\]
\item 
\[
\left\{\begin{aligned}
x-y+3z+4w+u-3v&=-18\\
2x+y-3z-2w+3u+3v&=11\\
2x-2y+6z+8w+2u-6v&=-35\\
\end{aligned}\right.
\]
\end{multicols}
\end{enumerate}
\end{Ejer}
%<<>>==<<>><<>>==<<>><<>>==<<>><<>>==<<>><<>>==<<>>
\section{Apendice: Factorizaciones}
%<<>>==<<>><<>>==<<>><<>>==<<>><<>>==<<>><<>>==<<>>
Este apéndice reúne un banco de ejemplos de tres factorizaciones matriciales: la factorización \(LU\), la factorización de Cholesky \(A=LL^{T}\) y la diagonalización \(A=PDP^{-1}\). Para cada método se incluyen cinco ejemplos con matrices \(2\times2\), cinco con matrices \(3\times3\) y cinco con matrices \(4\times4\), ordenados para disponer de varias alternativas de trabajo en clase, estudio y tarea.

En la factorización \(LU\) se trabaja sin pivoteo y sin normalizar pivotes. En Cholesky se emplean matrices simétricas definidas positivas. En la parte de diagonalización se inicia con la idea geométrica de direcciones que se conservan bajo la transformación y posteriormente se formaliza la descomposición \(A=PDP^{-1}\).

%<<>>==<<>><<>>==<<>><<>>==<<>><<>>==<<>><<>>==<<>>
\subsection*{Factorización $LU$}
%<<>>==<<>><<>>==<<>><<>>==<<>><<>>==<<>><<>>==<<>>
\subsubsection*{Recordatorio: de Gauss a $LU$}

Si no se requieren intercambios de renglones, hacemos eliminación de Gauss sin normalizar pivotes. Para anular $a_{ij}$ usamos
\[
m_{ij}=\frac{a_{ij}}{a_{jj}},\qquad R_i\leftarrow R_i-m_{ij}R_j.
\]
Los multiplicadores se guardan debajo de la diagonal de $L$ y la matriz final es $U$:
\[
\boxed{A=LU}.
\]

\begin{Ejer} LU $2\times 2$

Factorice, sin pivoteo,
\[
A=\begin{pmatrix}2 & -1\\-4 & 3\end{pmatrix}.
\]

\textbf{Objetivo}: Realizar Gauss sin normalizar los pivotes, guardar los multiplicadores y construir $L$ y $U$.
 Los multiplicadores de la eliminación son
\[
m_{21}=-2.
\]
Por tanto,
\[
L=\begin{pmatrix}1 & 0\\-2 & 1\end{pmatrix},\qquad U=\begin{pmatrix}2 & -1\\0 & 1\end{pmatrix}.
\]
Y se verifica directamente que
\[
\boxed{A=LU}.
\]
\end{Ejer}
%<<>>==<<>><<>>==<<>><<>>==<<>><<>>==<<>><<>>==<<>>
\begin{Ejer} 

Factorice, sin pivoteo,
\[
A=\begin{pmatrix}1 & 2\\3 & 9\end{pmatrix}.
\]
 Los multiplicadores de la eliminación son
\[
m_{21}=3.
\]
Por tanto,
\[
L=\begin{pmatrix}1 & 0\\3 & 1\end{pmatrix},\qquad U=\begin{pmatrix}1 & 2\\0 & 3\end{pmatrix}.
\]
Y se verifica directamente que
\[
\boxed{A=LU}.
\]

\end{Ejer}
%<<>>==<<>><<>>==<<>><<>>==<<>><<>>==<<>><<>>==<<>>
\begin{Ejer}
Factorice, sin pivoteo,
\[
A=\begin{pmatrix}3 & 0\\3 & -2\end{pmatrix}.
\]
 Los multiplicadores de la eliminación son
\[
m_{21}=1.
\]
Por tanto,
\[
L=\begin{pmatrix}1 & 0\\1 & 1\end{pmatrix},\qquad U=\begin{pmatrix}3 & 0\\0 & -2\end{pmatrix}.
\]
Y se verifica directamente que
\[
\boxed{A=LU}.
\]
\end{Ejer}
%<<>>==<<>><<>>==<<>><<>>==<<>><<>>==<<>><<>>==<<>>
\begin{Ejer}
Factorice, sin pivoteo,
\[
A=\begin{pmatrix}-2 & 3\\2 & 1\end{pmatrix}.
\]
Los multiplicadores de la eliminación son
\[
m_{21}=-1.
\]
Por tanto,
\[
L=\begin{pmatrix}1 & 0\\-1 & 1\end{pmatrix},\qquad U=\begin{pmatrix}-2 & 3\\0 & 4\end{pmatrix}.
\]
Y se verifica directamente que
\[
\boxed{A=LU}.
\]
\end{Ejer}
%<<>>==<<>><<>>==<<>><<>>==<<>><<>>==<<>><<>>==<<>>
\begin{Ejer} Factorice, sin pivoteo,
\[
A=\begin{pmatrix}4 & -2\\8 & -2\end{pmatrix}.
\]
 Los multiplicadores de la eliminación son
\[
m_{21}=2.
\]
Por tanto,
\[
L=\begin{pmatrix}1 & 0\\2 & 1\end{pmatrix},\qquad U=\begin{pmatrix}4 & -2\\0 & 2\end{pmatrix}.
\]
Y se verifica directamente que
\[
\boxed{A=LU}.
\]
\end{Ejer}
%<<>>==<<>><<>>==<<>><<>>==<<>><<>>==<<>><<>>==<<>>
\begin{Ejer}
Factorice, sin pivoteo,
\[
A=\begin{pmatrix}2 & -1 & 2\\-4 & 3 & -4\\-2 & 3 & 1\end{pmatrix}.
\]
Los multiplicadores de la eliminación son
\[
m_{21}=-2,\quad m_{31}=-1,\quad m_{32}=2.
\]
Por tanto,
\[
L=\begin{pmatrix}1 & 0 & 0\\-2 & 1 & 0\\-1 & 2 & 1\end{pmatrix},\qquad U=\begin{pmatrix}2 & -1 & 2\\0 & 1 & 0\\0 & 0 & 3\end{pmatrix}.
\]
Y se verifica directamente que
\[
\boxed{A=LU}.
\]

\end{Ejer}
%<<>>==<<>><<>>==<<>><<>>==<<>><<>>==<<>><<>>==<<>>
\begin{Ejer} Factorice, sin pivoteo,
\[
A=\begin{pmatrix}1 & 2 & 0\\3 & 9 & 3\\2 & -2 & -8\end{pmatrix}.
\]

Los multiplicadores de la eliminación son
\[
m_{21}=3,\quad m_{31}=2,\quad m_{32}=-2.
\]
Por tanto,
\[
L=\begin{pmatrix}1 & 0 & 0\\3 & 1 & 0\\2 & -2 & 1\end{pmatrix},\qquad U=\begin{pmatrix}1 & 2 & 0\\0 & 3 & 3\\0 & 0 & -2\end{pmatrix}.
\]
Y se verifica directamente que
\[
\boxed{A=LU}.
\]
\end{Ejer}
%<<>>==<<>><<>>==<<>><<>>==<<>><<>>==<<>><<>>==<<>>
\begin{Ejer}
Factorice, sin pivoteo,
\[
A=\begin{pmatrix}3 & 0 & 3\\3 & -2 & 1\\-6 & -6 & -8\end{pmatrix}.
\]

 Los multiplicadores de la eliminación son
\[
m_{21}=1,\quad m_{31}=-2,\quad m_{32}=3.
\]
Por tanto,
\[
L=\begin{pmatrix}1 & 0 & 0\\1 & 1 & 0\\-2 & 3 & 1\end{pmatrix},\qquad U=\begin{pmatrix}3 & 0 & 3\\0 & -2 & -2\\0 & 0 & 4\end{pmatrix}.
\]
Y se verifica directamente que
\[
\boxed{A=LU}.
\]
\end{Ejer}
%<<>>==<<>><<>>==<<>><<>>==<<>><<>>==<<>><<>>==<<>>
\begin{Ejer}
Factorice, sin pivoteo,
\[
A=\begin{pmatrix}-2 & 3 & -2\\2 & 1 & 3\\-6 & 13 & -3\end{pmatrix}.
\]


Los multiplicadores de la eliminación son
\[
m_{21}=-1,\quad m_{31}=3,\quad m_{32}=1.
\]
Por tanto,
\[
L=\begin{pmatrix}1 & 0 & 0\\-1 & 1 & 0\\3 & 1 & 1\end{pmatrix},\qquad U=\begin{pmatrix}-2 & 3 & -2\\0 & 4 & 1\\0 & 0 & 2\end{pmatrix}.
\]
Y se verifica directamente que
\[
\boxed{A=LU}.
\]
\end{Ejer}
%<<>>==<<>><<>>==<<>><<>>==<<>><<>>==<<>><<>>==<<>>
\begin{Ejer} Factorice, sin pivoteo,
\[
A=\begin{pmatrix}4 & -2 & 1\\8 & -2 & 1\\4 & -4 & 3\end{pmatrix}.
\]

Los multiplicadores de la eliminación son
\[
m_{21}=2,\quad m_{31}=1,\quad m_{32}=-1.
\]
Por tanto,
\[
L=\begin{pmatrix}1 & 0 & 0\\2 & 1 & 0\\1 & -1 & 1\end{pmatrix},\qquad U=\begin{pmatrix}4 & -2 & 1\\0 & 2 & -1\\0 & 0 & 1\end{pmatrix}.
\]
Y se verifica directamente que
\[
\boxed{A=LU}.
\]
\end{Ejer}
%<<>>==<<>><<>>==<<>><<>>==<<>><<>>==<<>><<>>==<<>>
\begin{Ejer}

Factorice, sin pivoteo,
\[
A=\begin{pmatrix}2 & -1 & 2 & 0\\-4 & 3 & -4 & 3\\-2 & 3 & 1 & 4\\6 & -2 & 3 & 3\end{pmatrix}.
\]

Los multiplicadores de la eliminación son
\[
m_{21}=-2,\quad m_{31}=-1,\quad m_{32}=2,\quad m_{41}=3,\quad m_{42}=1,\quad m_{43}=-1.
\]
Por tanto,
\[
L=\begin{pmatrix}1 & 0 & 0 & 0\\-2 & 1 & 0 & 0\\-1 & 2 & 1 & 0\\3 & 1 & -1 & 1\end{pmatrix},\qquad U=\begin{pmatrix}2 & -1 & 2 & 0\\0 & 1 & 0 & 3\\0 & 0 & 3 & -2\\0 & 0 & 0 & -2\end{pmatrix}.
\]
Y se verifica directamente que
\[
\boxed{A=LU}.
\]
\end{Ejer}
%<<>>==<<>><<>>==<<>><<>>==<<>><<>>==<<>><<>>==<<>>
\begin{Ejer}
Factorice, sin pivoteo,
\[
A=\begin{pmatrix}1 & 2 & 0 & 3\\3 & 9 & 3 & 7\\2 & -2 & -8 & 11\\1 & -1 & -7 & 11\end{pmatrix}.
\]
Los multiplicadores de la eliminación son
\[
m_{21}=3,\quad m_{31}=2,\quad m_{32}=-2,\quad m_{41}=1,\quad m_{42}=-1,\quad m_{43}=2.
\]
Por tanto,
\[
L=\begin{pmatrix}1 & 0 & 0 & 0\\3 & 1 & 0 & 0\\2 & -2 & 1 & 0\\1 & -1 & 2 & 1\end{pmatrix},\qquad U=\begin{pmatrix}1 & 2 & 0 & 3\\0 & 3 & 3 & -2\\0 & 0 & -2 & 1\\0 & 0 & 0 & 4\end{pmatrix}.
\]
Y se verifica directamente que
\[
\boxed{A=LU}.
\]
\end{Ejer}
%<<>>==<<>><<>>==<<>><<>>==<<>><<>>==<<>><<>>==<<>>
\begin{Ejer} Factorice, sin pivoteo,
\[
A=\begin{pmatrix}3 & 0 & 3 & -2\\3 & -2 & 1 & -1\\-6 & -6 & -8 & 6\\-3 & -4 & -15 & 8\end{pmatrix}.
\]

Los multiplicadores de la eliminación son
\[
m_{21}=1,\quad m_{31}=-2,\quad m_{32}=3,\quad m_{41}=-1,\quad m_{42}=2,\quad m_{43}=-2.
\]
Por tanto,
\[
L=\begin{pmatrix}1 & 0 & 0 & 0\\1 & 1 & 0 & 0\\-2 & 3 & 1 & 0\\-1 & 2 & -2 & 1\end{pmatrix},\qquad U=\begin{pmatrix}3 & 0 & 3 & -2\\0 & -2 & -2 & 1\\0 & 0 & 4 & -1\\0 & 0 & 0 & 2\end{pmatrix}.
\]
Y se verifica directamente que
\[
\boxed{A=LU}.
\]

\end{Ejer}
%<<>>==<<>><<>>==<<>><<>>==<<>><<>>==<<>><<>>==<<>>
\begin{Ejer} Factorice, sin pivoteo,
\[
A=\begin{pmatrix}-2 & 3 & -2 & 1\\2 & 1 & 3 & -2\\-6 & 13 & -3 & 4\\-4 & -2 & 0 & 11\end{pmatrix}.
\]

Los multiplicadores de la eliminación son
\[
m_{21}=-1,\quad m_{31}=3,\quad m_{32}=1,\quad m_{41}=2,\quad m_{42}=-2,\quad m_{43}=3.
\]
Por tanto,
\[
L=\begin{pmatrix}1 & 0 & 0 & 0\\-1 & 1 & 0 & 0\\3 & 1 & 1 & 0\\2 & -2 & 3 & 1\end{pmatrix},\qquad U=\begin{pmatrix}-2 & 3 & -2 & 1\\0 & 4 & 1 & -1\\0 & 0 & 2 & 2\\0 & 0 & 0 & 1\end{pmatrix}.
\]
Y se verifica directamente que
\[
\boxed{A=LU}.
\]
\end{Ejer}
%<<>>==<<>><<>>==<<>><<>>==<<>><<>>==<<>><<>>==<<>>
\begin{Ejer} Factorice, sin pivoteo,
\[
A=\begin{pmatrix}4 & -2 & 1 & -1\\8 & -2 & 1 & 0\\4 & -4 & 3 & -3\\-8 & 10 & -4 & 11\end{pmatrix}.
\]

Los multiplicadores de la eliminación son
\[
m_{21}=2,\quad m_{31}=1,\quad m_{32}=-1,\quad m_{41}=-2,\quad m_{42}=3,\quad m_{43}=1.
\]
Por tanto,
\[
L=\begin{pmatrix}1 & 0 & 0 & 0\\2 & 1 & 0 & 0\\1 & -1 & 1 & 0\\-2 & 3 & 1 & 1\end{pmatrix},\qquad U=\begin{pmatrix}4 & -2 & 1 & -1\\0 & 2 & -1 & 2\\0 & 0 & 1 & 0\\0 & 0 & 0 & 3\end{pmatrix}.
\]
Y se verifica directamente que
\[
\boxed{A=LU}.
\]
\end{Ejer}

%<<>>==<<>><<>>==<<>><<>>==<<>><<>>==<<>><<>>==<<>>
\section{Factorización de Cholesky}
%<<>>==<<>><<>>==<<>><<>>==<<>><<>>==<<>><<>>==<<>>

Para una matriz real simétrica definida positiva existe una matriz triangular inferior $L$, con diagonal positiva, tal que
\[
\boxed{A=LL^T}.
\]
Los elementos pueden obtenerse sucesivamente mediante
\[
l_{ii}=\sqrt{a_{ii}-\sum_{k<i}l_{ik}^2},\qquad
l_{ji}=\frac{a_{ji}-\sum_{k<i}l_{jk}l_{ik}}{l_{ii}}.
\]

\begin{Ejem}
Encuentre $L$ con diagonal positiva tal que
\[
A=\begin{pmatrix}1 & -1\\-1 & 5\end{pmatrix}=LL^T.
\]
Antes de calcular, observe que $A=A^T$.

Calculando los elementos de $L$ de izquierda a derecha y de arriba hacia abajo se obtiene
\[
l_{11}=1,\quad l_{21}=-1,\quad l_{22}=2.
\]
Así,
\[
L=\begin{pmatrix}1 & 0\\-1 & 2\end{pmatrix},\qquad L^T=\begin{pmatrix}1 & -1\\0 & 2\end{pmatrix}.
\]
Finalmente,
\[
\boxed{LL^T=A}.
\]
\end{Ejem}

\begin{Ejem}
Encuentre $L$ con diagonal positiva tal que
\[
A=\begin{pmatrix}4 & 4\\4 & 13\end{pmatrix}=LL^T.
\]
Antes de calcular, observe que $A=A^T$.

Calculando los elementos de $L$ de izquierda a derecha y de arriba hacia abajo se obtiene
\[
l_{11}=2,\quad l_{21}=2,\quad l_{22}=3.
\]
Así,
\[
L=\begin{pmatrix}2 & 0\\2 & 3\end{pmatrix},\qquad L^T=\begin{pmatrix}2 & 2\\0 & 3\end{pmatrix}.
\]
Finalmente,
\[
\boxed{LL^T=A}.
\]

\end{Ejem}

\begin{Ejem}

Encuentre $L$ con diagonal positiva tal que
\[
A=\begin{pmatrix}9 & -6\\-6 & 8\end{pmatrix}=LL^T.
\]
Antes de calcular, observe que $A=A^T$.

Calculando los elementos de $L$ de izquierda a derecha y de arriba hacia abajo se obtiene
\[
l_{11}=3,\quad l_{21}=-2,\quad l_{22}=2.
\]
Así,
\[
L=\begin{pmatrix}3 & 0\\-2 & 2\end{pmatrix},\qquad L^T=\begin{pmatrix}3 & -2\\0 & 2\end{pmatrix}.
\]
Finalmente,
\[
\boxed{LL^T=A}.
\]

\end{Ejem}

\begin{Ejem}

Encuentre $L$ con diagonal positiva tal que
\[
A=\begin{pmatrix}4 & 2\\2 & 2\end{pmatrix}=LL^T.
\]
Antes de calcular, observe que $A=A^T$.

Calculando los elementos de $L$ de izquierda a derecha y de arriba hacia abajo se obtiene
\[
l_{11}=2,\quad l_{21}=1,\quad l_{22}=1.
\]
Así,
\[
L=\begin{pmatrix}2 & 0\\1 & 1\end{pmatrix},\qquad L^T=\begin{pmatrix}2 & 1\\0 & 1\end{pmatrix}.
\]
Finalmente,
\[
\boxed{LL^T=A}.
\]

\end{Ejem}

\begin{Ejem}

Encuentre $L$ con diagonal positiva tal que
\[
A=\begin{pmatrix}4 & 6\\6 & 10\end{pmatrix}=LL^T.
\]
Antes de calcular, observe que $A=A^T$.

Calculando los elementos de $L$ de izquierda a derecha y de arriba hacia abajo se obtiene
\[
l_{11}=2,\quad l_{21}=3,\quad l_{22}=1.
\]
Así,
\[
L=\begin{pmatrix}2 & 0\\3 & 1\end{pmatrix},\qquad L^T=\begin{pmatrix}2 & 3\\0 & 1\end{pmatrix}.
\]
Finalmente,
\[
\boxed{LL^T=A}.
\]
\end{Ejem}

\begin{Ejem}

Encuentre $L$ con diagonal positiva tal que
\[
A=\begin{pmatrix}1 & -1 & -2\\-1 & 5 & 4\\-2 & 4 & 14\end{pmatrix}=LL^T.
\]
Antes de calcular, observe que $A=A^T$.

Calculando los elementos de $L$ de izquierda a derecha y de arriba hacia abajo se obtiene
\[
l_{11}=1,\quad l_{21}=-1,\quad l_{22}=2,\quad l_{31}=-2,\quad l_{32}=1,\quad l_{33}=3.
\]
Así,
\[
L=\begin{pmatrix}1 & 0 & 0\\-1 & 2 & 0\\-2 & 1 & 3\end{pmatrix},\qquad L^T=\begin{pmatrix}1 & -1 & -2\\0 & 2 & 1\\0 & 0 & 3\end{pmatrix}.
\]
Finalmente,
\[
\boxed{LL^T=A}.
\]

\end{Ejem}

\begin{Ejem}

Encuentre $L$ con diagonal positiva tal que
\[
A=\begin{pmatrix}4 & 4 & 2\\4 & 13 & 2\\2 & 2 & 5\end{pmatrix}=LL^T.
\]
Antes de calcular, observe que $A=A^T$.

Calculando los elementos de $L$ de izquierda a derecha y de arriba hacia abajo se obtiene
\[
l_{11}=2,\quad l_{21}=2,\quad l_{22}=3,\quad l_{31}=1,\quad l_{32}=0,\quad l_{33}=2.
\]
Así,
\[
L=\begin{pmatrix}2 & 0 & 0\\2 & 3 & 0\\1 & 0 & 2\end{pmatrix},\qquad L^T=\begin{pmatrix}2 & 2 & 1\\0 & 3 & 0\\0 & 0 & 2\end{pmatrix}.
\]
Finalmente,
\[
\boxed{LL^T=A}.
\]

\end{Ejem}

\begin{Ejem}

Encuentre $L$ con diagonal positiva tal que
\[
A=\begin{pmatrix}9 & -6 & 0\\-6 & 8 & -2\\0 & -2 & 2\end{pmatrix}=LL^T.
\]
Antes de calcular, observe que $A=A^T$.

Calculando los elementos de $L$ de izquierda a derecha y de arriba hacia abajo se obtiene
\[
l_{11}=3,\quad l_{21}=-2,\quad l_{22}=2,\quad l_{31}=0,\quad l_{32}=-1,\quad l_{33}=1.
\]
Así,
\[
L=\begin{pmatrix}3 & 0 & 0\\-2 & 2 & 0\\0 & -1 & 1\end{pmatrix},\qquad L^T=\begin{pmatrix}3 & -2 & 0\\0 & 2 & -1\\0 & 0 & 1\end{pmatrix}.
\]
Finalmente,
\[
\boxed{LL^T=A}.
\]
\end{Ejem}

\begin{Ejem}
Encuentre $L$ con diagonal positiva tal que
\[
A=\begin{pmatrix}4 & 2 & -2\\2 & 2 & 1\\-2 & 1 & 6\end{pmatrix}=LL^T.
\]
Antes de calcular, observe que $A=A^T$.


Calculando los elementos de $L$ de izquierda a derecha y de arriba hacia abajo se obtiene
\[
l_{11}=2,\quad l_{21}=1,\quad l_{22}=1,\quad l_{31}=-1,\quad l_{32}=2,\quad l_{33}=1.
\]
Así,
\[
L=\begin{pmatrix}2 & 0 & 0\\1 & 1 & 0\\-1 & 2 & 1\end{pmatrix},\qquad L^T=\begin{pmatrix}2 & 1 & -1\\0 & 1 & 2\\0 & 0 & 1\end{pmatrix}.
\]
Finalmente,
\[
\boxed{LL^T=A}.
\]

\end{Ejem}

\begin{Ejem}

Encuentre $L$ con diagonal positiva tal que
\[
A=\begin{pmatrix}1 & 0 & 2\\0 & 1 & -2\\2 & -2 & 12\end{pmatrix}=LL^T.
\]
Antes de calcular, observe que $A=A^T$.

Calculando los elementos de $L$ de izquierda a derecha y de arriba hacia abajo se obtiene
\[
l_{11}=1,\quad l_{21}=0,\quad l_{22}=1,\quad l_{31}=2,\quad l_{32}=-2,\quad l_{33}=2.
\]
Así,
\[
L=\begin{pmatrix}1 & 0 & 0\\0 & 1 & 0\\2 & -2 & 2\end{pmatrix},\qquad L^T=\begin{pmatrix}1 & 0 & 2\\0 & 1 & -2\\0 & 0 & 2\end{pmatrix}.
\]
Finalmente,
\[
\boxed{LL^T=A}.
\]
\end{Ejem}

\begin{Ejem}

Encuentre $L$ con diagonal positiva tal que
\[
A=\begin{pmatrix}1 & -1 & -2 & 0\\-1 & 5 & 4 & -2\\-2 & 4 & 14 & 5\\0 & -2 & 5 & 9\end{pmatrix}=LL^T.
\]
Antes de calcular, observe que $A=A^T$.

Calculando los elementos de $L$ de izquierda a derecha y de arriba hacia abajo se obtiene
\[
l_{11}=1,\quad l_{21}=-1,\quad l_{22}=2,\quad l_{31}=-2,\quad l_{32}=1,\quad l_{33}=3,\quad l_{41}=0,\quad l_{42}=-1,\quad l_{43}=2,\quad l_{44}=2.
\]
Así,
\[
L=\begin{pmatrix}1 & 0 & 0 & 0\\-1 & 2 & 0 & 0\\-2 & 1 & 3 & 0\\0 & -1 & 2 & 2\end{pmatrix},\qquad L^T=\begin{pmatrix}1 & -1 & -2 & 0\\0 & 2 & 1 & -1\\0 & 0 & 3 & 2\\0 & 0 & 0 & 2\end{pmatrix}.
\]
Finalmente,
\[
\boxed{LL^T=A}.
\]


\end{Ejem}

\begin{Ejem}

Encuentre $L$ con diagonal positiva tal que
\[
A=\begin{pmatrix}4 & 4 & 2 & -2\\4 & 13 & 2 & 4\\2 & 2 & 5 & -5\\-2 & 4 & -5 & 10\end{pmatrix}=LL^T.
\]
Antes de calcular, observe que $A=A^T$.

Calculando los elementos de $L$ de izquierda a derecha y de arriba hacia abajo se obtiene
\[
l_{11}=2,\quad l_{21}=2,\quad l_{22}=3,\quad l_{31}=1,\quad l_{32}=0,\quad l_{33}=2,\quad l_{41}=-1,\quad l_{42}=2,\quad l_{43}=-2,\quad l_{44}=1.
\]
Así,
\[
L=\begin{pmatrix}2 & 0 & 0 & 0\\2 & 3 & 0 & 0\\1 & 0 & 2 & 0\\-1 & 2 & -2 & 1\end{pmatrix},\qquad L^T=\begin{pmatrix}2 & 2 & 1 & -1\\0 & 3 & 0 & 2\\0 & 0 & 2 & -2\\0 & 0 & 0 & 1\end{pmatrix}.
\]
Finalmente,
\[
\boxed{LL^T=A}.
\]

\end{Ejem}

\begin{Ejem}
Encuentre $L$ con diagonal positiva tal que
\[
A=\begin{pmatrix}9 & -6 & 0 & 6\\-6 & 8 & -2 & -8\\0 & -2 & 2 & 3\\6 & -8 & 3 & 10\end{pmatrix}=LL^T.
\]
Antes de calcular, observe que $A=A^T$.

Calculando los elementos de $L$ de izquierda a derecha y de arriba hacia abajo se obtiene
\[
l_{11}=3,\quad l_{21}=-2,\quad l_{22}=2,\quad l_{31}=0,\quad l_{32}=-1,\quad l_{33}=1,\quad l_{41}=2,\quad l_{42}=-2,\quad l_{43}=1,\quad l_{44}=1.
\]
Así,
\[
L=\begin{pmatrix}3 & 0 & 0 & 0\\-2 & 2 & 0 & 0\\0 & -1 & 1 & 0\\2 & -2 & 1 & 1\end{pmatrix},\qquad L^T=\begin{pmatrix}3 & -2 & 0 & 2\\0 & 2 & -1 & -2\\0 & 0 & 1 & 1\\0 & 0 & 0 & 1\end{pmatrix}.
\]
Finalmente,
\[
\boxed{LL^T=A}.
\]

\end{Ejem}

\begin{Ejem}
Encuentre $L$ con diagonal positiva tal que
\[
A=\begin{pmatrix}4 & 2 & -2 & -4\\2 & 2 & 1 & -1\\-2 & 1 & 6 & 4\\-4 & -1 & 4 & 9\end{pmatrix}=LL^T.
\]
Antes de calcular, observe que $A=A^T$.

Calculando los elementos de $L$ de izquierda a derecha y de arriba hacia abajo se obtiene
\[
l_{11}=2,\quad l_{21}=1,\quad l_{22}=1,\quad l_{31}=-1,\quad l_{32}=2,\quad l_{33}=1,\quad l_{41}=-2,\quad l_{42}=1,\quad l_{43}=0,\quad l_{44}=2.
\]
Así,
\[
L=\begin{pmatrix}2 & 0 & 0 & 0\\1 & 1 & 0 & 0\\-1 & 2 & 1 & 0\\-2 & 1 & 0 & 2\end{pmatrix},\qquad L^T=\begin{pmatrix}2 & 1 & -1 & -2\\0 & 1 & 2 & 1\\0 & 0 & 1 & 0\\0 & 0 & 0 & 2\end{pmatrix}.
\]
Finalmente,
\[
\boxed{LL^T=A}.
\]

\end{Ejem}

\begin{Ejem}

Encuentre $L$ con diagonal positiva tal que
\[
A=\begin{pmatrix}1 & 0 & 2 & 1\\0 & 1 & -2 & 0\\2 & -2 & 12 & 0\\1 & 0 & 0 & 11\end{pmatrix}=LL^T.
\]
Antes de calcular, observe que $A=A^T$.


Calculando los elementos de $L$ de izquierda a derecha y de arriba hacia abajo se obtiene
\[
l_{11}=1,\quad l_{21}=0,\quad l_{22}=1,\quad l_{31}=2,\quad l_{32}=-2,\quad l_{33}=2,\quad l_{41}=1,\quad l_{42}=0,\quad l_{43}=-1,\quad l_{44}=3.
\]
Así,
\[
L=\begin{pmatrix}1 & 0 & 0 & 0\\0 & 1 & 0 & 0\\2 & -2 & 2 & 0\\1 & 0 & -1 & 3\end{pmatrix},\qquad L^T=\begin{pmatrix}1 & 0 & 2 & 1\\0 & 1 & -2 & 0\\0 & 0 & 2 & -1\\0 & 0 & 0 & 3\end{pmatrix}.
\]
Finalmente,
\[
\boxed{LL^T=A}.
\]
\end{Ejem}

%<<>>==<<>>==<<>>==<<>>==<<>>==<<>>==<<>>==<<>>==<<>>==
\section{Diagonalización}
%<<>>==<<>>==<<>>==<<>>==<<>>==<<>>==<<>>==<<>>==<<>>==

Buscamos vectores no nulos cuya dirección no cambie al aplicar $A$:
\[
A\mathbf v=\lambda\mathbf v.
\]
En esas direcciones la transformación sólo multiplica por un número. Si encontramos suficientes direcciones independientes, podemos usarlas como una nueva base.
\medskip
En esa base, la acción de $A$ queda descrita por una matriz diagonal. Ésta es la idea que conduce a
\[
\boxed{A=PDP^{-1}}.
\]
Las columnas de $P$ son vectores propios y la diagonal de $D$ contiene sus valores propios en el mismo orden.

\begin{Ejem}

Diagonalice, si es posible,
\[
A=\begin{pmatrix}1 & 0\\2 & 3\end{pmatrix}.
\]

Busque direcciones $v$ para las cuales $Av=\lambda v$ y compruebe que se obtienen $2$ vectores propios linealmente independientes.

\scriptsize
Una elección de pares propios es
\[
\lambda_1=1,\;v_1=\begin{pmatrix}1\\-1\end{pmatrix};\quad \lambda_2=3,\;v_2=\begin{pmatrix}0\\1\end{pmatrix}.
\]
Con estos vectores,
\[
P=\begin{pmatrix}1 & 0\\-1 & 1\end{pmatrix},\qquad D=\begin{pmatrix}1 & 0\\0 & 3\end{pmatrix},\qquad P^{-1}=\begin{pmatrix}1 & 0\\1 & 1\end{pmatrix}.
\]
Por construcción $AP=PD$; equivalentemente,
\[
\boxed{A=PDP^{-1}}.
\]
\end{Ejem}

\begin{Ejem}

Diagonalice, si es posible,
\[
A=\begin{pmatrix}2 & 0\\6 & -1\end{pmatrix}.
\]


Busque direcciones $v$ para las cuales $Av=\lambda v$ y compruebe que se obtienen $2$ vectores propios linealmente independientes.

\scriptsize
Una elección de pares propios es
\[
\lambda_1=2,\;v_1=\begin{pmatrix}1\\2\end{pmatrix};\quad \lambda_2=-1,\;v_2=\begin{pmatrix}0\\1\end{pmatrix}.
\]
Con estos vectores,
\[
P=\begin{pmatrix}1 & 0\\2 & 1\end{pmatrix},\qquad D=\begin{pmatrix}2 & 0\\0 & -1\end{pmatrix},\qquad P^{-1}=\begin{pmatrix}1 & 0\\-2 & 1\end{pmatrix}.
\]
Por construcción $AP=PD$; equivalentemente,
\[
\boxed{A=PDP^{-1}}.
\]

\end{Ejem}

\begin{Ejem}


Diagonalice, si es posible,
\[
A=\begin{pmatrix}4 & 0\\-6 & 1\end{pmatrix}.
\]

\scriptsize
Una elección de pares propios es
\[
\lambda_1=4,\;v_1=\begin{pmatrix}1\\-2\end{pmatrix};\quad \lambda_2=1,\;v_2=\begin{pmatrix}0\\1\end{pmatrix}.
\]
Con estos vectores,
\[
P=\begin{pmatrix}1 & 0\\-2 & 1\end{pmatrix},\qquad D=\begin{pmatrix}4 & 0\\0 & 1\end{pmatrix},\qquad P^{-1}=\begin{pmatrix}1 & 0\\2 & 1\end{pmatrix}.
\]
Por construcción $AP=PD$; equivalentemente,
\[
\boxed{A=PDP^{-1}}.
\]
\end{Ejem}

\begin{Ejem}


Diagonalice, si es posible,
\[
A=\begin{pmatrix}-2 & 0\\-5 & 3\end{pmatrix}.
\]

\scriptsize
Una elección de pares propios es
\[
\lambda_1=-2,\;v_1=\begin{pmatrix}1\\1\end{pmatrix};\quad \lambda_2=3,\;v_2=\begin{pmatrix}0\\1\end{pmatrix}.
\]
Con estos vectores,
\[
P=\begin{pmatrix}1 & 0\\1 & 1\end{pmatrix},\qquad D=\begin{pmatrix}-2 & 0\\0 & 3\end{pmatrix},\qquad P^{-1}=\begin{pmatrix}1 & 0\\-1 & 1\end{pmatrix}.
\]
Por construcción $AP=PD$; equivalentemente,
\[
\boxed{A=PDP^{-1}}.
\]

\end{Ejem}

\begin{Ejem}

Diagonalice, si es posible,
\[
A=\begin{pmatrix}5 & 0\\-3 & 2\end{pmatrix}.
\]


\scriptsize
Una elección de pares propios es
\[
\lambda_1=5,\;v_1=\begin{pmatrix}1\\-1\end{pmatrix};\quad \lambda_2=2,\;v_2=\begin{pmatrix}0\\1\end{pmatrix}.
\]
Con estos vectores,
\[
P=\begin{pmatrix}1 & 0\\-1 & 1\end{pmatrix},\qquad D=\begin{pmatrix}5 & 0\\0 & 2\end{pmatrix},\qquad P^{-1}=\begin{pmatrix}1 & 0\\1 & 1\end{pmatrix}.
\]
Por construcción $AP=PD$; equivalentemente,
\[
\boxed{A=PDP^{-1}}.
\]

\end{Ejem}

\begin{Ejem}

Diagonalice, si es posible,
\[
A=\begin{pmatrix}1 & 0 & 0\\1 & 2 & 0\\-8 & -2 & 4\end{pmatrix}.
\]

\scriptsize
Una elección de pares propios es
\[
\lambda_1=1,\;v_1=\begin{pmatrix}1\\-1\\2\end{pmatrix};\quad \lambda_2=2,\;v_2=\begin{pmatrix}0\\1\\1\end{pmatrix};\quad \lambda_3=4,\;v_3=\begin{pmatrix}0\\0\\1\end{pmatrix}.
\]
Con estos vectores,
\[
P=\begin{pmatrix}1 & 0 & 0\\-1 & 1 & 0\\2 & 1 & 1\end{pmatrix},\qquad D=\begin{pmatrix}1 & 0 & 0\\0 & 2 & 0\\0 & 0 & 4\end{pmatrix},\qquad P^{-1}=\begin{pmatrix}1 & 0 & 0\\1 & 1 & 0\\-3 & -1 & 1\end{pmatrix}.
\]
Por construcción $AP=PD$; equivalentemente,
\[
\boxed{A=PDP^{-1}}.
\]

\end{Ejem}

\begin{Ejem}


Diagonalice, si es posible,
\[
A=\begin{pmatrix}3 & 0 & 0\\8 & -1 & 0\\-6 & 3 & 2\end{pmatrix}.
\]

\paragraph{Guía}
Busque direcciones $v$ para las cuales $Av=\lambda v$ y compruebe que se obtienen $3$ vectores propios linealmente independientes.

\scriptsize
Una elección de pares propios es
\[
\lambda_1=3,\;v_1=\begin{pmatrix}1\\2\\0\end{pmatrix};\quad \lambda_2=-1,\;v_2=\begin{pmatrix}0\\1\\-1\end{pmatrix};\quad \lambda_3=2,\;v_3=\begin{pmatrix}0\\0\\1\end{pmatrix}.
\]
Con estos vectores,
\[
P=\begin{pmatrix}1 & 0 & 0\\2 & 1 & 0\\0 & -1 & 1\end{pmatrix},\qquad D=\begin{pmatrix}3 & 0 & 0\\0 & -1 & 0\\0 & 0 & 2\end{pmatrix},\qquad P^{-1}=\begin{pmatrix}1 & 0 & 0\\-2 & 1 & 0\\-2 & 1 & 1\end{pmatrix}.
\]
Por construcción $AP=PD$; equivalentemente,
\[
\boxed{A=PDP^{-1}}.
\]

\end{Ejem}

\begin{Ejem}

Diagonalice, si es posible,
\[
A=\begin{pmatrix}2 & 0 & 0\\6 & 5 & 0\\32 & 14 & -2\end{pmatrix}.
\]
\scriptsize
Una elección de pares propios es
\[
\lambda_1=2,\;v_1=\begin{pmatrix}1\\-2\\1\end{pmatrix};\quad \lambda_2=5,\;v_2=\begin{pmatrix}0\\1\\2\end{pmatrix};\quad \lambda_3=-2,\;v_3=\begin{pmatrix}0\\0\\1\end{pmatrix}.
\]
Con estos vectores,
\[
P=\begin{pmatrix}1 & 0 & 0\\-2 & 1 & 0\\1 & 2 & 1\end{pmatrix},\qquad D=\begin{pmatrix}2 & 0 & 0\\0 & 5 & 0\\0 & 0 & -2\end{pmatrix},\qquad P^{-1}=\begin{pmatrix}1 & 0 & 0\\2 & 1 & 0\\-5 & -2 & 1\end{pmatrix}.
\]
Por construcción $AP=PD$; equivalentemente,
\[
\boxed{A=PDP^{-1}}.
\]
\end{Ejem}

\begin{Ejem}


Diagonalice, si es posible,
\[
A=\begin{pmatrix}-1 & 0 & 0\\-4 & 3 & 0\\-2 & 2 & 4\end{pmatrix}.
\]

\scriptsize
Una elección de pares propios es
\[
\lambda_1=-1,\;v_1=\begin{pmatrix}1\\1\\0\end{pmatrix};\quad \lambda_2=3,\;v_2=\begin{pmatrix}0\\1\\-2\end{pmatrix};\quad \lambda_3=4,\;v_3=\begin{pmatrix}0\\0\\1\end{pmatrix}.
\]
Con estos vectores,
\[
P=\begin{pmatrix}1 & 0 & 0\\1 & 1 & 0\\0 & -2 & 1\end{pmatrix},\qquad D=\begin{pmatrix}-1 & 0 & 0\\0 & 3 & 0\\0 & 0 & 4\end{pmatrix},\qquad P^{-1}=\begin{pmatrix}1 & 0 & 0\\-1 & 1 & 0\\-2 & 2 & 1\end{pmatrix}.
\]
Por construcción $AP=PD$; equivalentemente,
\[
\boxed{A=PDP^{-1}}.
\]

\end{Ejem}

\begin{Ejem}

Diagonalice, si es posible,
\[
A=\begin{pmatrix}4 & 0 & 0\\-3 & 1 & 0\\-9 & -5 & 6\end{pmatrix}.
\]

\paragraph{Guía}
Busque direcciones $v$ para las cuales $Av=\lambda v$ y compruebe que se obtienen $3$ vectores propios linealmente independientes.

\scriptsize
Una elección de pares propios es
\[
\lambda_1=4,\;v_1=\begin{pmatrix}1\\-1\\2\end{pmatrix};\quad \lambda_2=1,\;v_2=\begin{pmatrix}0\\1\\1\end{pmatrix};\quad \lambda_3=6,\;v_3=\begin{pmatrix}0\\0\\1\end{pmatrix}.
\]
Con estos vectores,
\[
P=\begin{pmatrix}1 & 0 & 0\\-1 & 1 & 0\\2 & 1 & 1\end{pmatrix},\qquad D=\begin{pmatrix}4 & 0 & 0\\0 & 1 & 0\\0 & 0 & 6\end{pmatrix},\qquad P^{-1}=\begin{pmatrix}1 & 0 & 0\\1 & 1 & 0\\-3 & -1 & 1\end{pmatrix}.
\]
Por construcción $AP=PD$; equivalentemente,
\[
\boxed{A=PDP^{-1}}.
\]
\end{Ejem}

\begin{Ejem}


Diagonalice, si es posible,
\[
A=\begin{pmatrix}1 & 0 & 0 & 0\\1 & 2 & 0 & 0\\-5 & -1 & 3 & 0\\-9 & -1 & 4 & 5\end{pmatrix}.
\]

\scriptsize
Una elección de pares propios es
\[
\lambda_1=1,\;v_1=\begin{pmatrix}1\\-1\\2\\0\end{pmatrix};\quad \lambda_2=2,\;v_2=\begin{pmatrix}0\\1\\1\\-1\end{pmatrix};\quad \lambda_3=3,\;v_3=\begin{pmatrix}0\\0\\1\\-2\end{pmatrix};\quad \lambda_4=5,\;v_4=\begin{pmatrix}0\\0\\0\\1\end{pmatrix}.
\]
Con estos vectores,
\[
P=\begin{pmatrix}1 & 0 & 0 & 0\\-1 & 1 & 0 & 0\\2 & 1 & 1 & 0\\0 & -1 & -2 & 1\end{pmatrix},\qquad D=\begin{pmatrix}1 & 0 & 0 & 0\\0 & 2 & 0 & 0\\0 & 0 & 3 & 0\\0 & 0 & 0 & 5\end{pmatrix},\qquad P^{-1}=\begin{pmatrix}1 & 0 & 0 & 0\\1 & 1 & 0 & 0\\-3 & -1 & 1 & 0\\-5 & -1 & 2 & 1\end{pmatrix}.
\]
Por construcción $AP=PD$; equivalentemente,
\[
\boxed{A=PDP^{-1}}.
\]

\end{Ejem}

\begin{Ejem}


Diagonalice, si es posible,
\[
A=\begin{pmatrix}2 & 0 & 0 & 0\\6 & -1 & 0 & 0\\-10 & 5 & 4 & 0\\-3 & 1 & 1 & 3\end{pmatrix}.
\]
\scriptsize
Una elección de pares propios es
\[
\lambda_1=2,\;v_1=\begin{pmatrix}1\\2\\0\\1\end{pmatrix};\quad \lambda_2=-1,\;v_2=\begin{pmatrix}0\\1\\-1\\0\end{pmatrix};\quad \lambda_3=4,\;v_3=\begin{pmatrix}0\\0\\1\\1\end{pmatrix};\quad \lambda_4=3,\;v_4=\begin{pmatrix}0\\0\\0\\1\end{pmatrix}.
\]
Con estos vectores,
\[
P=\begin{pmatrix}1 & 0 & 0 & 0\\2 & 1 & 0 & 0\\0 & -1 & 1 & 0\\1 & 0 & 1 & 1\end{pmatrix},\qquad D=\begin{pmatrix}2 & 0 & 0 & 0\\0 & -1 & 0 & 0\\0 & 0 & 4 & 0\\0 & 0 & 0 & 3\end{pmatrix},\qquad P^{-1}=\begin{pmatrix}1 & 0 & 0 & 0\\-2 & 1 & 0 & 0\\-2 & 1 & 1 & 0\\1 & -1 & -1 & 1\end{pmatrix}.
\]
Por construcción $AP=PD$; equivalentemente,
\[
\boxed{A=PDP^{-1}}.
\]
\end{Ejem}

\begin{Ejem}


Diagonalice, si es posible,
\[
A=\begin{pmatrix}3 & 0 & 0 & 0\\-4 & 1 & 0 & 0\\17 & 6 & -2 & 0\\-19 & -6 & 7 & 5\end{pmatrix}.
\]

\scriptsize
Una elección de pares propios es
\[
\lambda_1=3,\;v_1=\begin{pmatrix}1\\-2\\1\\0\end{pmatrix};\quad \lambda_2=1,\;v_2=\begin{pmatrix}0\\1\\2\\-2\end{pmatrix};\quad \lambda_3=-2,\;v_3=\begin{pmatrix}0\\0\\1\\-1\end{pmatrix};\quad \lambda_4=5,\;v_4=\begin{pmatrix}0\\0\\0\\1\end{pmatrix}.
\]
Con estos vectores,
\[
P=\begin{pmatrix}1 & 0 & 0 & 0\\-2 & 1 & 0 & 0\\1 & 2 & 1 & 0\\0 & -2 & -1 & 1\end{pmatrix},\qquad D=\begin{pmatrix}3 & 0 & 0 & 0\\0 & 1 & 0 & 0\\0 & 0 & -2 & 0\\0 & 0 & 0 & 5\end{pmatrix},\qquad P^{-1}=\begin{pmatrix}1 & 0 & 0 & 0\\2 & 1 & 0 & 0\\-5 & -2 & 1 & 0\\-1 & 0 & 1 & 1\end{pmatrix}.
\]
Por construcción $AP=PD$; equivalentemente,
\[
\boxed{A=PDP^{-1}}.
\]

\end{Ejem}

\begin{Ejem}


Diagonalice, si es posible,
\[
A=\begin{pmatrix}-2 & 0 & 0 & 0\\-3 & 1 & 0 & 0\\-6 & 6 & 4 & 0\\-8 & -8 & -4 & 6\end{pmatrix}.
\]

\scriptsize
Una elección de pares propios es
\[
\lambda_1=-2,\;v_1=\begin{pmatrix}1\\1\\0\\2\end{pmatrix};\quad \lambda_2=1,\;v_2=\begin{pmatrix}0\\1\\-2\\0\end{pmatrix};\quad \lambda_3=4,\;v_3=\begin{pmatrix}0\\0\\1\\2\end{pmatrix};\quad \lambda_4=6,\;v_4=\begin{pmatrix}0\\0\\0\\1\end{pmatrix}.
\]
Con estos vectores,
\[
P=\begin{pmatrix}1 & 0 & 0 & 0\\1 & 1 & 0 & 0\\0 & -2 & 1 & 0\\2 & 0 & 2 & 1\end{pmatrix},\qquad D=\begin{pmatrix}-2 & 0 & 0 & 0\\0 & 1 & 0 & 0\\0 & 0 & 4 & 0\\0 & 0 & 0 & 6\end{pmatrix},\qquad P^{-1}=\begin{pmatrix}1 & 0 & 0 & 0\\-1 & 1 & 0 & 0\\-2 & 2 & 1 & 0\\2 & -4 & -2 & 1\end{pmatrix}.
\]
Por construcción $AP=PD$; equivalentemente,
\[
\boxed{A=PDP^{-1}}.
\]

\end{Ejem}

\begin{Ejem}

Diagonalice, si es posible,
\[
A=\begin{pmatrix}5 & 0 & 0 & 0\\-3 & 2 & 0 & 0\\15 & 3 & -1 & 0\\-23 & -7 & 8 & 3\end{pmatrix}.
\]

\scriptsize
Una elección de pares propios es
\[
\lambda_1=5,\;v_1=\begin{pmatrix}1\\-1\\2\\0\end{pmatrix};\quad \lambda_2=2,\;v_2=\begin{pmatrix}0\\1\\1\\-1\end{pmatrix};\quad \lambda_3=-1,\;v_3=\begin{pmatrix}0\\0\\1\\-2\end{pmatrix};\quad \lambda_4=3,\;v_4=\begin{pmatrix}0\\0\\0\\1\end{pmatrix}.
\]
Con estos vectores,
\[
P=\begin{pmatrix}1 & 0 & 0 & 0\\-1 & 1 & 0 & 0\\2 & 1 & 1 & 0\\0 & -1 & -2 & 1\end{pmatrix},\qquad D=\begin{pmatrix}5 & 0 & 0 & 0\\0 & 2 & 0 & 0\\0 & 0 & -1 & 0\\0 & 0 & 0 & 3\end{pmatrix},\qquad P^{-1}=\begin{pmatrix}1 & 0 & 0 & 0\\1 & 1 & 0 & 0\\-3 & -1 & 1 & 0\\-5 & -1 & 2 & 1\end{pmatrix}.
\]
Por construcción $AP=PD$; equivalentemente,
\[
\boxed{A=PDP^{-1}}.
\]
\end{Ejem}
\end{document}
